\documentclass[10pt,letterpaper,twocolumn]{article}
\usepackage[margin=0.68in,columnsep=0.24in]{geometry}
\usepackage[T1]{fontenc}
\usepackage{mathptmx}
\usepackage{microtype,amsmath,amssymb,amsthm,mathtools}
\usepackage{graphicx,booktabs,array,enumitem,xcolor,caption,url,balance}
\usepackage[colorlinks=true,allcolors=blue!45!black]{hyperref}
\captionsetup{font=footnotesize,labelfont=bf}
\setlength{\parindent}{1em}
\setlength{\parskip}{1pt}
\raggedbottom
\setlist{nosep,leftmargin=*}
\newtheorem{theorem}{Theorem}
\newtheorem{lemma}[theorem]{Lemma}
\newtheorem{proposition}[theorem]{Proposition}
\newtheorem{corollary}[theorem]{Corollary}
\newcommand{\F}{\mathbb F_2}
\newcommand{\rank}{\operatorname{rank}_{\F}}
\newcommand{\supp}{\operatorname{supp}}
\newcommand{\own}{\operatorname{own}}
\newcommand{\im}{\operatorname{im}}
\newcommand{\Span}{\operatorname{span}}
\newcommand{\cost}{\mathcal C}
\newcommand{\reach}{\mathcal V}
\newcommand{\xor}{\oplus}
\newcommand{\AlgBox}[2]{\par\smallskip\noindent\fbox{\begin{minipage}{0.95\columnwidth}\small\textbf{#1}\par\smallskip #2\end{minipage}}\par\smallskip}
\renewcommand{\thesection}{\Roman{section}}
\renewcommand{\thesubsection}{\Alph{subsection}}
\makeatletter
\renewcommand{\p@subsection}{\thesection-}
\makeatother
\title{\vspace{-1.5em}Seamwise: Boundary Repair for Partitioned\\Decoding of Logical Quantum Circuits}
\author{Songhao Zhang\textsuperscript{1}\qquad Tianjie Hu\textsuperscript{1}\\[3pt]\small\textsuperscript{1}Taiyi Quantum Science \& Technology Co., Ltd.}
\date{}
\begin{document}
\maketitle
\begin{abstract}
Transversal entangling gates spread physical faults across encoded blocks, creating detector evidence that a partitioned decoder can separate prematurely. Seamwise repairs these boundaries by reopening nearby fault variables, conditioning on the fixed exterior, and accepting only syndrome-preserving improvements in configuration likelihood. We connect the logical freedom of these repairs to their algebraic execution cost. A rank identity characterizes the logical changes available to each repair band, while Seamwise's merge-persistent hypergraph union--find oracle retains disjoint detector-supported bases across cluster merges and eliminates each newly internal fault column once during growth. For repair calls containing $W$ fault-column incidences in total and maximum Tanner-cluster size $K_*$, a support-local implementation has deterministic repair work $O(IW+WK_*^2+W\log(W+2))$ under bounded detector degree and the stated growth scheduler, where $I$ bounds the BP iterations. A dependent-column certificate distinguishes full-band logical freedom from the logical directions exposed by growth. Circuit experiments evaluate the outer repair policy and find 30.6--34.5\% fewer logical failures than direct cutting in the primary compact surface-code comparison. On 256 identical distance-five circuit samples, a BP-assisted reference implementation returns identical corrections with 29 logical failures while reducing mean full repair latency from 143.4 to 132.0 ms when its rebuilding oracle is replaced by the persistent oracle. Native Tesseract yields ten failures at 169.9 ms on the same samples, exposing an accuracy--latency trade-off. We additionally measure process memory and scale block count, code distance, and circuit depth. The framework explicitly couples repair locality, accessible logical transformations, and the algebraic work required to execute them.
\end{abstract}

\section{Introduction}
Transversal entangling gates couple error correction across encoded blocks. A physical fault on one block can propagate through a gate and flip detectors on several blocks, as the path in Fig.~\ref{fig:overview}(a) illustrates. The resulting detector outcomes provide evidence about the same underlying fault. A joint decoder can use that evidence in a single circuit-wide inference problem~\cite{cain}. Partitioning makes each inference problem smaller, but a partition boundary can separate the detector evidence needed to distinguish two plausible fault configurations (Fig.~\ref{fig:overview}(b)).

Our direct-cut baseline processes partitions in block order. After committing each local correction, it forwards that correction's detector effects to later partitions through a residual syndrome. Forwarding communicates the consequences of a choice, but the later partition cannot revise the fault variables that produced it. This distinction matters because two configurations can reproduce the same measured syndrome while implying different corrections to the logical outputs. A cut correction can therefore satisfy every detector and still make a logical error.

We introduce Seamwise, a refinement procedure that starts from a syndrome-consistent correction. At each selected boundary, or seam, it reopens a band of fault variables assigned to blocks on either side and fixes the correction everywhere else. It removes the exterior contribution from every detector equation touched by the band, including equations whose detectors lie beyond the band. An inner decoder then proposes a local replacement. Seamwise accepts a proposal only if it satisfies all of those equations and increases the prior likelihood of the fault configuration. The incumbent remains a feasible solution throughout.

The repair band controls which logical alternatives the decoder can reach. We derive a rank identity that counts the independent changes to monitored logical outputs attainable while preserving the syndrome and fixing the exterior. In the four-block distance-five surface-code and $[[72,12,6]]$ bivariate bicycle (BB) circuits, radius-one bands contain about half of the fault variables and span the full monitored logical-change space. The acceptance rule improves configuration likelihood. We measure logical accuracy separately because logical-class probabilities aggregate over multiple configurations. We also show that repairs on disjoint fault-variable sets commute even when their detector footprints overlap, allowing concurrent updates with the same correction as serial execution under the stated solver conditions.

Repair locality alone does not control the cost of repeatedly solving the local systems. Seamwise therefore maintains a merge-persistent HUF state inside each conditional solve. Its state consists of a binary basis, a reduced target residual, and a correction witness. When Tanner components merge, their old bases occupy disjoint detector coordinates and concatenate without cross-elimination. Only newly internal fault columns require elimination. This provides a cluster-sensitive work bound for the repair stage without changing ownership, conditioning, or the likelihood acceptance rule. The bound applies to this explicit backend, not automatically to all four previously measured inner decoders.

We evaluate the repair policy on matched fault samples across four inner-decoder settings. In the main benchmark, with four $[[25,1,5]]$ surface-code blocks and circuit noise probability $p=0.006$, radius-one repair yields 30.6--34.5\% fewer logical failures than direct cutting on 256 paired shots. All four comparisons are significant after Holm adjustment across the 135 planned tests. The BB experiments test the method on high-rate blocks under phenomenological noise. On eight distance-five surface-code blocks, parallel repair reduces mean warm-call latency from 60.00 to 49.71 ms relative to serial repair, while direct cutting takes 35.44 ms. The controlled backend study additionally measures rebuilding and persistent HUF on identical stored circuit samples, with native Tesseract and stronger-order BP postprocessing as external comparisons.

The paper makes three contributions:
\begin{enumerate}
\item \textbf{A syndrome-preserving local repair decoder.} Ownership, residual forwarding, and conditioned seam replacement form a complete algorithm with proved consistency and strict likelihood descent.
\item \textbf{A connection between logical reach and algebraic work.} The original rank identity characterizes a band's available logical transformations; the integrated persistent oracle supplies certified local corrections with a deterministic cluster-sensitive work bound. A dependent-column certificate identifies which logical directions the growth trajectory actually exposes.
\item \textbf{Paired accuracy, execution, and scale measurements.} Circuit comparisons measure the repair policy, exact-oracle ablations isolate persistent computation, and native external baselines expose the remaining accuracy--latency gap.
\end{enumerate}

\begin{figure*}[t]
\centering\includegraphics[width=\textwidth]{figures/fig1.pdf}
\caption{Fault propagation and partition boundaries. (a) A physical fault can propagate across encoded blocks. (b) The dashed lines mark boundaries between schematic decoding subproblems. In the implementation, partitions are defined by fault ownership across encoded blocks. Seamwise reopens a band of fault variables at a boundary and retains every detector constraint that depends on the band.}\label{fig:overview}
\end{figure*}

\section{Related work}
Correlated decoding uses fault propagation through transversal gates to improve logical-circuit inference~\cite{cain}. Cain et al. subsequently formulate efficient decoding in terms of propagating logical-operator products~\cite{fast}. Our method operates directly on a fixed detector model and repairs boundaries introduced by its spatial partition.

Parallel window decoding establishes overlapping windows and correction reconciliation as scalable decoding tools~\cite{window}. Boundary repair complements pre-commit overlap with a post-cut operation: the existing correction supplies a feasible incumbent, and the exterior remains fixed while the seam is refined.

BB codes provide high-rate compact quantum LDPC blocks~\cite{bb,bbpre}; the bicycle architecture develops transversal logical instructions and Bell-reference simulation methods~\cite{tour}. We combine those Bell-reference observables with explicit surface-code circuits and two-sector BB phenomenological circuits. We evaluate three postprocessors for belief propagation (BP). BeliefFind (BF) uses union-find inversion. The other two use localized statistics decoding (LSD)~\cite{lsd} and ordered statistics decoding (OSD)~\cite{osd}.

Localized statistics decoding develops on-the-fly elimination and merge-aware local factorization~\cite{lsd}. Seamwise specializes its persistent state to HUF's full internal-hyperedge eligibility rule~\cite{cain} within fixed-exterior post-cut repair. Its analysis connects exact stopping and witness certificates with logical-change spaces and the work of a complete repair schedule.

\section{Detector model and decoding objective}
\subsection{Model and notation}
Table~\ref{tab:notation} collects the notation. We use $[[n_q,k_q,d]]$ for a single code block and reserve $N_f,N_{\rm det},N_{\rm obs}$ for circuit fault mechanisms, detector bits, and monitored logical-output bits.

Let $D\in\F^{N_{\rm det}\times N_f}$ be the detector matrix and $L\in\F^{N_{\rm obs}\times N_f}$ the logical-action matrix. Column $j$ represents a binary fault mechanism, with detector signature $D_{:,j}$ and logical signature $L_{:,j}$. For a sampled fault vector $e\in\F^{N_f}$,
\begin{equation}s=De,\qquad \ell=Le.\label{eq:model}\end{equation}
All matrix operations on syndromes and fault vectors are over $\F$. A decoder returns an estimate $\hat e$ satisfying $D\hat e=s$, or declares failure. A syndrome-consistent estimate succeeds logically precisely when
\begin{equation}L(\hat e+e)=0.\label{eq:logical}\end{equation}
The experimental failure statistic is the whole-circuit probability that any monitored logical output is incorrect.

For independent fault probabilities $p_j\in(0,1)$, the configuration objective, up to an additive constant, is
\begin{equation}\cost(x)=\sum_{j=1}^{N_f}w_jx_j,\qquad w_j=\log\frac{1-p_j}{p_j}.\label{eq:cost}\end{equation}
The prior weights in Eq.~\eqref{eq:cost} are fixed throughout repair. BP posterior reliabilities may guide the inner search, but acceptance uses the original configuration objective. We evaluate logical performance using Eq.~\eqref{eq:logical} and accept repairs using the configuration objective. Appendix~\ref{app:examples} gives the corresponding logical-class formulation.

\begin{table}[t]\centering\small
\caption{Core notation for the detector model, repair, and persistent backend.}\label{tab:notation}
\begin{tabular}{@{}p{0.24\columnwidth}p{0.69\columnwidth}@{}}\toprule
Symbol&Meaning\\\midrule
$D,L$&Maps from faults to detector bits and logical flips.\\
$e,\hat e$&Sampled fault vector and decoded estimate.\\
$s,\ell$&Measured syndrome $De$ and sampled logical flips $Le$.\\
$x,z$&Current estimate and candidate seam replacement.\\
$J,\bar J,R(J)$&Reopened columns, fixed exterior, and all touched detector rows.\\
$D_J,t_J$&$D_{R(J),J}$ and its conditioned target.\\
$w_j,\cost$&Prior fault weight and configuration cost.\\
$\rho(J),\reach(J)$&Dimension and subspace of available logical changes.\\
$S,\mathcal I(S)$&Tanner component and its eligible internal columns.\\
$\mathcal B_S,r_S,u_S$&Persistent basis, target residual, correction witness.\\
$m_a,K_a$&Fault count and maximum cluster size in repair call $a$.\\
$W,K_*$&$\sum_a m_a$ and $\max_aK_a$.\\
$F_a,\rho_{\rm act}$&Columns actually inserted and their exposed logical reach.\\\bottomrule
\end{tabular}\end{table}

\subsection{Constructing the matrices from a small circuit}
Prepare three qubits $a,b,c$ in $|000\rangle$, apply CNOTs $a\to b$ and $a\to c$, and finally measure all three in $Z$ (Fig.~\ref{fig:small}). This is a repetition-code circuit restricted to $X$ faults. Write its measurement bits as $u_a,u_b,u_c$. The two detector bits are the stabilizer parities $s_1=u_a+u_b$ and $s_2=u_b+u_c$, corresponding to $Z_aZ_b$ and $Z_bZ_c$. The monitored logical observable is $\bar Z=Z_a$, so its flip bit is $\ell=u_a$. All three parities have known zero values in the ideal circuit.

Each column is constructed by activating one fault and recording the resulting detector and logical flips. Figure~\ref{fig:small} gives the circuit and the four fault signatures. Stacking these signatures as columns gives
\begin{equation}
D=\begin{pmatrix}1&1&0&1\\0&1&1&1\end{pmatrix},\qquad L=\begin{pmatrix}1&0&0&1\end{pmatrix}.\label{eq:toy}
\end{equation}
Here $N_f=4$, $N_{\rm det}=2$, and $N_{\rm obs}=1$: a column counts a fault mechanism at a circuit location, rather than a physical qubit. For example, $e=(0,0,0,1)^T$ produces $s=(1,1)^T$ and $\ell=1$. The alternative $\hat e=(0,1,0,0)^T$ has the same syndrome but $L\hat e=0$. Thus matching detector bits alone does not determine the logical action. In the full experiments, the same column construction includes measurement faults and propagated multi-block faults, with final Bell products providing the monitored logical observables.

\begin{figure*}[t]\centering\includegraphics[width=\textwidth]{figures/fig2.pdf}
\caption{Constructing detector and logical-action columns. Each $Xe_j$ marks an optional fault in the circuit (left). Activating that fault alone gives row $j$ of the signature table (right). Fault 4 occurs before the second CNOT, which propagates its $X_a$ to $X_aX_c$.}\label{fig:small}\end{figure*}

\section{Boundary-repair algorithm and guarantees}
\subsection{Fault ownership and residual syndromes}
Assign every detector row $i$ a block label $\beta(i)\in\{1,\ldots,B\}$. For each nonzero detector column, define its owner by
\begin{equation}\own(j)=\min\{\beta(i):D_{ij}=1\}.\label{eq:owner}\end{equation}
Ownership is determined by the propagated detector signature. It ensures that a fault owned by a later block has no detector effect in an earlier block:
\begin{equation}\beta(i)<\own(j)\quad\Longrightarrow\quad D_{ij}=0.\label{eq:triangular}\end{equation}
Syndrome-invisible columns are assigned their prior-minimizing values. Their sampled logical effects remain included in Eq.~\eqref{eq:logical}.

Let $F$ be the set of already committed variables, with committed estimate $x_F$. The residual syndrome is
\begin{equation}s^{\rm res}=s+D_{:,F}x_F.\label{eq:forward}\end{equation}
Forwarding means applying Eq.~\eqref{eq:forward} to all affected rows, including rows in later blocks. Updating only the rows of the current core would lose precisely the inter-block information that correlated decoding is intended to retain.

\subsection{Initial partition pass}
The initial pass decodes two-owner-block cores in order, commits each core's variables, and forwards their complete detector signature into the residual syndrome. A halo of width $h$ adds uncommitted owner blocks as context. Core consistency is checked at every commitment; the same-engine joint solver supplies the fallback. Appendix~\ref{app:initial} gives the full algorithm and the core-closure proof.

\subsection{Post-cut seam repair}
Begin with a syndrome-consistent direct-cut estimate $x$. A seam at block boundary $b$ separates blocks $b-1$ and $b$. For radius $r$, unfreeze the owner band
\begin{equation}J_{b,r}=\{j:b-r\le\own(j)<b+r\},\label{eq:band}\end{equation}
with block labels clipped to the valid range. Define its complete touched-row set
\begin{equation}R(J)=\{i:\exists j\in J,\ D_{ij}=1\}.\label{eq:rows}\end{equation}
The exterior $x_{\bar J}$ is fixed, so the replacement target is
\begin{equation}t_J=s_{R(J)}+D_{R(J),\bar J}x_{\bar J}.\label{eq:target}\end{equation}
Subscripts select detector rows first and fault columns second: $D_{R(J),J}$ describes the unfrozen variables, while $D_{R(J),\bar J}x_{\bar J}$ is the known detector contribution of the exterior. Appendix~\ref{app:initial} gives a numerical example.

Re-decode $D_{R(J),J}z=t_J$ and accept a candidate only when it is consistent and
\begin{equation}\sum_{j\in J}w_jz_j<\sum_{j\in J}w_jx_j.\label{eq:accept}\end{equation}
Each seam therefore defines a constrained refinement of the existing correction, rather than a new independent decoding problem.

\AlgBox{Algorithm 1: likelihood-filtered seam repair}{
\textbf{Input:} feasible $x$, ordered seams $\mathcal S$, radius $r$, priors, inner solver $\mathcal A$.
\begin{enumerate}
\item For the next seam, form $J$, $R(J)$, and $t_J$ from the current $x$.
\item Compute $z\leftarrow\mathcal A(D_{R(J),J},t_J,p_J)$.
\item If $D_{R(J),J}z=t_J$ and Eq.~\eqref{eq:accept} holds, set $x_J\leftarrow z$. Otherwise retain $x_J$.
\end{enumerate}
Repeat for each seam in order; return $x$.}

\begin{theorem}[Feasible descent]\label{thm:descent}
Every accepted replacement in Algorithm 1 preserves $Dx=s$ and strictly decreases $\cost(x)$. Any sequence of accepted replacements is finite, even if seams are revisited.
\end{theorem}
\begin{proof}
On $R(J)$, Eq.~\eqref{eq:target} enforces the original syndrome after the exterior is restored. Outside $R(J)$, every column in $J$ is zero, so replacing $x_J$ changes no detector. The exterior contribution to Eq.~\eqref{eq:cost} is fixed, and Eq.~\eqref{eq:accept} therefore strictly decreases total cost. There are finitely many binary configurations; strict descent prevents a configuration from recurring.
\end{proof}
The incumbent restriction $x_J$ is a feasible solution of the conditioned problem. Retaining it when a candidate is rejected makes each band update a complete, consistency-preserving decoder operation. A radius-one band comprises two owner blocks. Its detector footprint can extend farther because a fault owned in the band may propagate into later blocks. This is why all touched rows, rather than only rows whose labels lie inside the band, are required in Eq.~\eqref{eq:rows}.

\subsection{Which bands can change a logical decision?}
For a band $J$, write $D_J=D_{R(J),J}$ and define
\begin{equation}\rho(J)=\rank\begin{pmatrix}D_J\\L_{:,J}\end{pmatrix}-\rank(D_J).\label{eq:rank}\end{equation}
\begin{theorem}[Logical freedom of a repair band]\label{thm:reach}
Fix the exterior of a syndrome-consistent correction. The set of logical outputs attainable by feasible replacements on $J$ is an affine binary space of dimension $\rho(J)$. In particular, $\rho(J)=0$ if and only if every feasible replacement preserves the complete monitored logical output.
\end{theorem}
\begin{proof}
Any two feasible replacements differ by $v\in\ker D_J$, and every such $v$ gives another feasible replacement. Their logical differences form $\reach(J)=L_{:,J}(\ker D_J)$. By rank--nullity its dimension is
\begin{align*}
\dim(\ker D_J)&-\dim(\ker D_J\cap\ker L_{:,J})\\
&=\rank\begin{pmatrix}D_J\\L_{:,J}\end{pmatrix}-\rank(D_J).
\end{align*}
Translation by the incumbent logical output gives the claimed affine space.
\end{proof}
The identity turns repair width into a precise measure of logical reach. A zero-rank band preserves the monitored logical output; a rank-$\rho$ band admits $2^\rho$ feasible logical offsets relative to its incumbent. For nested bands $J\subseteq J'$, zero-extension embeds $\ker D_J$ into $\ker D_{J'}$, giving $\rho(J)\le\rho(J')$. Section~\ref{sec:results} measures this hierarchy in compact surface and BB blocks. Logical reach describes representability, not the success probability of a heuristic search or a minimum-likelihood guarantee.

\subsection{Conflict graph and simultaneous repair}
Let a complete feasible correction be $x$. A feasible update on a band $J$ has a difference vector $\delta_J$, supported on $J$, satisfying $D\delta_J=0$. Construct a conflict graph with one vertex per repair band and an edge exactly when two bands share a fault variable.
\begin{theorem}[Commuting disjoint-band repair]\label{thm:commute}
Suppose bands $J_1,\ldots,J_a$ are pairwise disjoint. Compute a candidate for each band from the same feasible incumbent $x$, using a solver whose candidate depends only on the fixed local matrix, target, and priors. The feasible, strictly improving replacements can be accepted simultaneously. Their combined correction equals the result of applying the same candidates in any serial order. Detector-row overlap does not require an additional conflict edge.
\end{theorem}
\begin{proof}
Every accepted difference lies in $\ker D$, so their sum also lies in $\ker D$. For $J\cap K=\varnothing$, a feasible update $\delta_K$ changes the conditional target on $J$ by $D_{R(J),K}(\delta_K)_K=0$. It changes neither the incumbent variables on $J$ nor their cost. Thus the candidate, its feasibility test, and its acceptance decision are unchanged. Disjoint replacements consequently commute, and their strict cost decreases add.
\end{proof}
A proper coloring of the variable-conflict graph defines an execution schedule: process colors in order and the bands of each color concurrently. Updates in different colors may interact, so the color order is part of the algorithm. With core width two, radius-one seam bands are disjoint and form one color; the radius-two bands use two colors in the tested chains (Fig.~\ref{fig:bands}). The implementation validates bitwise equivalence to serial radius-one repair.

For a single disjoint batch, a band with $\rho(J)=0$ can also be skipped without changing the batch's logical output: its update has zero logical action and leaves every other band's target and acceptance test unchanged. The rank-screened variant applies this rule to the disjoint, single-pass radius-one schedule.

\subsection{Merge-persistent conditional decoding}\label{sec:meld}
The outer repair policy is unchanged. For a selected band, put $A=D_J$, $t=t_J$, and $m=|J|$, omitting the already fixed syndrome-invisible columns. Because $Dx=s$,
\begin{equation}t_J=D_Jx_J.\label{eq:incumbenttarget}\end{equation}
Thus a repair target is feasible by construction. This identity also permits target construction by multiplying the incumbent's local columns, without scanning a potentially large exterior. It does not make every intermediate HUF cluster feasible.

In the Tanner graph of $A$, retain all touched detector rows and only the unfrozen fault columns. For a component $S$, let $C(S)$ be its check vertices and
\begin{equation}\mathcal I(S)=\{j:f_j\in S,\ \supp(A_{:,j})\subseteq C(S)\}.\label{eq:internal}\end{equation}
A column is legal only after its entire detector neighborhood is internal. Restricting an exterior-touching column to the current component would change the problem.

For each component, maintain pivot pairs $(b_p,c_p)$ with $Ac_p=b_p$, a residual $r_S$, and a witness $u_S$ such that
\begin{equation}
\begin{split}
t_S&=r_S+Au_S,\quad r_S[p]=0\ \text{at every stored pivot},\\
\Span\{b_p\}&=\im A_{:,\mathcal I(S)}.
\end{split}\label{eq:invariant}
\end{equation}
The symbols $r_S$ and $\rho(J)$ denote different quantities: the first is a syndrome vector, the second a logical-space dimension. Separate Tanner components have disjoint check coordinates. Their old bases therefore concatenate without cross-elimination; residuals and witnesses add. Each newly internal column is reduced once against the inherited basis, with coefficient vectors updated alongside it. The cluster is valid exactly when $r_S=0$. Appendix~\ref{app:persistent} gives the update algorithm and proof.

\begin{theorem}[Exact conditional oracle]\label{thm:oracle}
On a fixed local matrix, target, and merge-only HUF trajectory, the persistent oracle returns exactly the validity decisions of rebuilding every internal-column system. A valid witness obeys $Au_S=t_S$ and uses only internal columns. With identical growth weights, selection policy, and tie breaking, it induces the same growth trajectory as the rebuilding oracle.
\end{theorem}
Theorem~\ref{thm:oracle} preserves feasibility decisions. Distinct satisfying witnesses may have distinct logical actions and costs. When identical decoder outputs are required, one terminal canonical solve per final cluster reproduces the rebuilding implementation's fixed deterministic column-order/pivot rule within the bound below. Algorithm 1 still checks both the complete local syndrome and the original prior cost. A valid HUF witness is a proposal, not a minimum-cost solution.

\begin{corollary}[End-to-end correction preservation]\label{cor:outerequivalence}
Fix the model, priors, deterministic BP preprocessing, growth ties, terminal canonicalization, initial partition and fallback policies, and repair schedule. In executions that complete, replacing each rebuilding HUF oracle by its persistent counterpart leaves the entire Seamwise correction vector and every repair-acceptance decision unchanged. Consequently their logical-failure indicators agree pointwise on every sampled fault vector.
\end{corollary}
\begin{proof}
For each conditional call, Theorem~\ref{thm:oracle} and terminal canonicalization give the same validity flag and candidate. Induction through ownership commitments and fallbacks gives the same initial incumbent. Induction through the repair schedule then gives the same conditioned targets, BP inputs, candidates, prior-cost tests and updated incumbents. The common final vector has the same logical action against any true fault vector. External process deadlines are outside this equivalence statement.
\end{proof}

\subsection{Repair work and exposed logical reach}\label{sec:work}
Consider a finite repair schedule with calls indexed by $a$. Let $m_a=|J_a|$, $K_a$ be the maximum number of check and fault vertices in a component during call $a$, $W=\sum_a m_a$, and $K_*=\max_aK_a$. Repeated visits and overlapping bands contribute repeatedly to $W$.

\begin{theorem}[Cluster-sensitive repair work]\label{thm:work}
Assume uniformly bounded row and column degree of every local detector matrix; merge-only growth with nonnegative finite boundary weights; constant-cost, bounded-precision weight updates; and a scheduler with $O(m_a)$ successful growth events, $O(m_aK_a)$ incidence scans/root queries, and $O(m_a)$ priority-queue updates per call. With support-local exact XOR algebra and at most $I$ BP iterations, the persistent repair backend has deterministic work
\begin{equation}
T_{\rm repair}=O\!\left(\sum_a\left[Im_a+m_aK_a^2+m_a\log(m_a+2)\right]\right),\label{eq:workexact}
\end{equation}
and hence
\begin{equation}\boxed{T_{\rm repair}=O\!\left(IW+WK_*^2+W\log(W+2)\right).}\label{eq:work}
\end{equation}
A sequential implementation uses $O(N_f+\max_a m_aK_a)$ space beyond the stored detector model and any offline rank tables. Optional terminal canonicalization is included in these bounds.
\end{theorem}
All complexity derivations, including the scheduling assumptions, are in Appendix~\ref{app:complexity}. If every nonzero fault column is visited at most $\mu$ times, then $W\le\mu N_f$. For bounded $I,\mu$ and $K_*=O(\log N_f)$, Eq.~\eqref{eq:work} gives $O(N_f\log^2N_f)$ repair work. The unconditional worst case permits $K_a=O(m_a)$ and gives $O(\sum_a m_a^3)$ algebraic work. This is a guarantee for the repair stage with this backend; the initial partition pass, global fallbacks, preprocessing, and any extra candidate search must be added separately. A fixed schedule bounds $W$; finite descent alone does not bound the number of revisits polynomially.

The rank identity in Theorem~\ref{thm:reach} concerns every column in $J$. HUF can terminate after exposing only $F\subseteq J$. During persistent insertion, a dependent column reduces to zero with a coefficient relation $k_j$ satisfying $A k_j=0$. Retain its logical image $\eta_j=L_{:,J}k_j$ instead of discarding that relation entirely.

\begin{proposition}[Exposed logical-reach certificate]\label{prop:exposed}
Let $F$ contain all columns inserted during a local solve. The dependent-column relations form a basis of $\ker D_{:,F}$ after zero extension. Consequently
\begin{equation}
\begin{split}
\reach_{\rm act}(F)&=\Span\{\eta_j:j\text{ inserted dependently}\}\\
&=L_{:,F}(\ker D_{:,F})\subseteq\reach(J),\\
\rho_{\rm act}(F)&=\dim\reach_{\rm act}(F)\le\rho(J).
\end{split}\label{eq:active}
\end{equation}
Thus $\rho_{\rm act}(F)<\rho(J)$ certifies that the inserted support cannot realize all logical offsets available to the full band. Equality certifies equal logical-change spaces, not equal cost optima or logical-class probabilities.
\end{proposition}
This certificate connects the two scales of locality. Seamwise chooses which fault variables may be reconsidered; persistent HUF chooses which of those variables enter its algebra. In Eq.~\eqref{eq:toy}, insert columns 1 and 2 first: the inserted matrix has no kernel, so $\rho_{\rm act}=0$. Inserting column 4 then produces $k_4=(0,1,0,1)^T$, with $Dk_4=0$ and $Lk_4=1$. The exposed logical rank becomes one without rebuilding the detector basis. Appendix~\ref{app:reach} proves the certificate and accounts for its additional cost. No logical-rank guarantee is inferred from cluster size alone.

The certificate is optional instrumentation. It does not change the reported policy or retrospectively enlarge the original circuit study. A future policy could continue growth until a prescribed reach is exposed, but its altered stopping rule, cost, and logical performance would require a separate evaluation.


\begin{figure*}[t]\centering\includegraphics[width=\textwidth]{figures/fig3.pdf}
\caption{Variable supports for eight owner blocks and core width two. A halo adds context before committing the core. Radius-one repair uses three disjoint bands in one color; radius-two bands overlap and require two colors. Rectangles describe fault-owner support; the full touched-row detector footprint is defined by Eq.~\eqref{eq:rows}.}\label{fig:bands}\end{figure*}

\section{Experimental design}
\paragraph{Compact code blocks.}
The main experiments use rotated $[[9,1,3]]$ and $[[25,1,5]]$ surface-code blocks and $[[72,12,6]]$ BB blocks. Four-block circuits monitor eight logical Bell outputs for the surface code and 96 for BB. An eight-block distance-five surface circuit monitors 16 outputs and exercises multiple repair seams. Every block begins in an ideal encoded Bell pair with noiseless references, followed by transversal CNOT matchings, noisy syndrome extraction, and ideal final closure~\cite{tour}. The decoder receives the complete syndrome record and partitions it by spatial fault ownership. Surface-code extraction uses explicit ancilla schedules with depolarizing gate noise, reset flips, and readout flips at probability $p$. The BB control uses data depolarization and readout-flip noise with ideal transversal CNOTs and ideal stabilizer-product measurements. Stim generates the exact undecomposed detector models~\cite{stim}. Appendix~\ref{app:protocol} gives all circuit configurations and implementation details.

\paragraph{Matched engine and policy comparisons.}
BF, BP+LSD0, BP+LSD2, and BP+OSD0 each run joint decoding, direct cutting, and radius-one repair on identical fault samples~\cite{bfdocs,lsddocs,osddocs}. They share 30 min-sum BP iterations, scaling 0.625, parallel message scheduling, and one native CPU thread. BF additionally tests halo, radius-two, and rank-screened variants. A radius-one band contains two owner blocks; every replacement conditions on its exterior and retains the complete touched-row detector constraints. In four-block circuits, radius two covers all nonzero detector columns; in eight-block circuits, its bands contain four owner blocks.

\paragraph{Paired inference and execution measurements.}
The complete original study contains 4,096 seeded circuit samples, 15 configurations, and 16 decoder stacks. The frozen protocol specifies nine contrasts per configuration. We use exact conditional McNemar tests~\cite{mcnemar} with Holm adjustment across all 135 contrasts~\cite{holm}. The plots show 95\% marginal Wilson intervals~\cite{wilson}. The full outcomes, effect intervals, seeds, and protocol history are retained in the original study. Engineering measurements use three fresh-process replicates on 64 common inputs, exclude startup from warm-call latency, and separately record startup and process-tree memory. Reused timing and functional-audit inputs do not increase the independent logical sample count.

\paragraph{Controlled backend and external-baseline study.}
The new study reuses stored circuits, detector matrices, and noise samples from the original repair-policy grid. It compares a rebuilding HUF reference with the persistent conditional backend inside the same ownership/cut/repair wrapper. Both use identical weighted boundary growth and terminal canonicalization. Tesseract~\cite{tesseract} and native BP+LSD2/OSD2 are evaluated on the same complete detector records. Appendix~\ref{app:study} specifies the pilot budgets, scaling circuits, timing scopes, and unavailable baselines.

\section{Compact-code results}\label{sec:results}
\subsection{Logical failure rates across inner decoders}
The primary benchmark uses four $[[25,1,5]]$ blocks at $p=0.006$. Radius-one boundary repair reduces BF failures from 65 to 45 in 256 shared shots, a 30.8\% relative reduction. The paired outcomes contain 21 corrected failures and one new failure, giving a Holm-adjusted $p=0.00147$. BP+LSD2 improves from 58 to 38 failures, a 34.5\% reduction with adjusted $p=0.00477$. All four cut-to-repair contrasts in Table~\ref{tab:primary} are significant after Holm adjustment across the 135 planned comparisons.

\begin{table*}[t]\centering\small
\caption{Logical failures before and after radius-one repair on 256 paired shots for four $[[25,1,5]]$ blocks at $p=0.006$. Paired gain is the absolute decrease in failure rate in percentage points (pp), with a conservative 95\% interval.}\label{tab:primary}
\begin{tabular}{lrrrrrr}\toprule
Engine&Cut&Repair&Fixed&New&Paired gain, pp [95\% CI]&Holm $p$\\\midrule
BF&65&45&21&1&7.8 [2.3, 12.9]&0.00147\\
LSD0&65&43&24&2&8.6 [2.6, 14.2]&0.00142\\
LSD2&58&38&22&2&7.8 [2.0, 13.3]&0.00477\\
OSD0&62&43&21&2&7.4 [1.7, 12.8]&0.00872\\\bottomrule
\end{tabular}\end{table*}
In the $d=5$, $p=0.003$ case, BF failures decrease from 55 to 39 on 1,024 paired shots. Figure~\ref{fig:failures} shows the surface-code and BB72 operating points with Wilson intervals.

\begin{figure*}[t]\centering\includegraphics[width=\textwidth]{figures/fig4.pdf}
\caption{BF direct cutting and radius-one repair on paired samples. The $d=3$ surface-code and BB72 cases use 512 and 128 shots per point, respectively. The $d=5$ cases use 1,024 shots at $p=0.003$ and 256 at $p=0.006$. Error bars show marginal 95\% Wilson intervals.}\label{fig:failures}\end{figure*}

\subsection{Extension to compact high-rate BB blocks}
The $[[72,12,6]]$ BB case extends boundary repair to blocks that each encode 12 logical qubits. At $p=0.015$, BP+LSD2 reduces the observed count from 44 to 30 failures in 128 shots, a 31.8\% reduction; BF reduces 44 to 32, a 27.3\% reduction. Each shot simultaneously monitors both Pauli sectors through 96 Bell-reference observables. These results describe the tested phenomenological BB noise model.

\subsection{Half-size bands with full logical reach}
The rank identity explains why a narrow owner band can address a rich set of circuit-level logical alternatives. In the four-block distance-five model, a radius-one band contains 5,299 of the full problem's 10,164 fault variables and retains all eight dimensions of monitored logical freedom. In BB72 it contains exactly half the fault variables and retains all 96 logical-change dimensions. The compact $d=3$ band accesses seven of eight dimensions. Table~\ref{tab:reach} makes this local-to-global relationship explicit.
\begin{table*}[t]\centering\small
\caption{Logical reach of a radius-one band in four-block circuits.}\label{tab:reach}
\begin{tabular}{lrrr}\toprule
Code block&Fault variables (local / full)&Share (\%)&$\rho$ (local / full)\\\midrule
Surface $[[9,1,3]]$&949 / 1,821&52.1&7 / 8\\
Surface $[[25,1,5]]$&5,299 / 10,164&52.1&8 / 8\\
BB $[[72,12,6]]$&3,456 / 6,912&50.0&96 / 96\\\bottomrule
\end{tabular}\end{table*}
Full logical reach in the distance-five and BB72 bands means that every monitored logical offset is attainable by a syndrome-consistent replacement with the exterior fixed. The inner decoder and acceptance rule determine which feasible replacements are made in practice. The nested-support measurements in Appendix~\ref{app:surface} resolve how that reach emerges as columns enter a band. Proposition~\ref{prop:exposed} additionally separates this full-band capacity from the support actually inserted by persistent HUF; the original rank table does not measure the latter.

\subsection{Parallel repair on eight compact surface-code blocks}
The eight-block $[[25,1,5]]$ circuit has three disjoint radius-one seam bands. Two worker processes execute those updates concurrently and return bitwise-identical corrections to the serial implementation. Mean warm-call latency decreases from 60.00 to 49.71 ms, corresponding to $1.21\times$ the serial-repair throughput. Direct cutting takes 35.44 ms. Parallel repair increases peak process-tree proportional set size (PSS) from 88.1 to 233.1 MiB relative to serial repair (Table~\ref{tab:parallel}). At distance three, parallel repair increases mean warm-call latency from 3.53 to 4.57 ms (Table~\ref{tab:engineering}). Warm-call timings pool three fresh-process replicates. PSS is sampled separately on the first eight inputs of each run. Appendix~\ref{app:engineering} gives the complete execution grid.

All six radius-one process-parallel trials reproduce serial corrections bit for bit. An additional 32-input functional audit executes overlapping radius-two bands in two colors and matches the serial reference with the same color order. The implementation therefore realizes both the disjoint-update composition theorem and a concrete conflict schedule.
\begin{table}[t]\centering\small
\caption{Warm-call performance and logical failures on 64 shared shots for eight $[[25,1,5]]$ blocks at $p=0.003$.}\label{tab:parallel}
\begin{tabular}{lrrrr}\toprule
Method&Failures&Mean ms&p95 ms&PSS MiB\\\midrule
Joint BF&9 / 64&68.78&149.46&74.6\\
Cut BF&9 / 64&35.44&77.20&82.3\\
Serial repair&3 / 64&60.00&137.54&88.1\\
Parallel repair&3 / 64&49.71&122.24&233.1\\\bottomrule
\end{tabular}\end{table}

\subsection{Controlled persistent-backend circuit pilot}\label{sec:newpilot}
Across 24 circuit configurations and 168 planned configuration/backend jobs, 136 jobs complete the full measurement protocol. The 37 matched rebuilding/persistent comparisons cover 2,560 paired shot executions, with zero correction-vector or recorded growth-history disagreements. Full details and non-completions are retained in the archive.
\begin{table*}[t]\centering\scriptsize
\caption{Completed controlled pilot runs. HUF/SW: complete Seamwise with prior-weighted rebuilding/persistent HUF. BPHUF/BPSW: matched BP-assisted counterparts. Tesseract, LSD2 and OSD2: native joint decoders. RSS is the separately sampled whole-process footprint, not backend-only storage.}\label{tab:newpilot}
\begin{tabular}{llrrrr}\toprule Configuration&Backend&Failures/shots&Mean ms&p95 ms&RSS MiB\\\midrule
bb144\_p0.005\_B4&BPHUF&0/64&134.74&209.08&136.4\\
bb144\_p0.005\_B4&BPSW&0/64&132.83&204.10&137.1\\
bb144\_p0.005\_B4&HUF&2/64&224.01&376.17&126.5\\
bb144\_p0.005\_B4&LSD2&0/64&13.46&21.00&74.7\\
bb144\_p0.005\_B4&OSD2&0/64&959.25&1151.50&175.5\\
bb144\_p0.005\_B4&SW&2/64&225.91&374.91&126.5\\
bb144\_p0.005\_B4&Tesseract&0/64&14.75&19.29&117.4\\
surface\_d3\_p0.003\_B4&BPHUF&2/64&1.99&6.91&75.6\\
surface\_d3\_p0.003\_B4&BPSW&2/64&2.05&6.99&75.4\\
surface\_d3\_p0.003\_B4&HUF&20/64&6.22&17.67&74.0\\
surface\_d3\_p0.003\_B4&LSD2&2/64&0.39&1.29&69.0\\
surface\_d3\_p0.003\_B4&OSD2&3/64&1.55&7.24&72.6\\
surface\_d3\_p0.003\_B4&SW&20/64&6.68&18.91&73.7\\
surface\_d3\_p0.003\_B4&Tesseract&2/64&2.15&8.39&97.1\\
surface\_d5\_p0.006\_B4&BPHUF&29/256&143.39&369.52&113.1\\
surface\_d5\_p0.006\_B4&BPSW&29/256&132.00&329.58&113.0\\
surface\_d5\_p0.006\_B4&LSD2&47/256&15.05&47.81&73.2\\
surface\_d5\_p0.006\_B4&Tesseract&10/256&169.88&432.11&113.7\\
\bottomrule\end{tabular}\end{table*}
\begin{figure*}[t]\centering\includegraphics[width=\textwidth]{study_scaling.pdf}\caption{Measured scaling pilot at fixed noise p=.003. Block sweep fixes d=3 and T=3; distance sweep fixes four blocks and T=3; depth sweep fixes d=3 and four blocks. Error bars are marginal Wilson intervals. Only completed, fixed-sample runs are plotted; fixed shallow extraction depth is not a threshold experiment.}\label{fig:newscale}\end{figure*}
\begin{figure*}[t]\centering\includegraphics[width=\textwidth]{study_clusters.pdf}\caption{Peak-cluster distributions over recorded completed solver calls. Each panel uses configurations completed by both members of that oracle pair. BP early exits have zero peak; infeasible rejected-call histories are unavailable. The two panels need not have the same configuration cohort.}\label{fig:newclusters}\end{figure*}
The controlled study does not establish a general latency advantage of the current persistent implementation. The exact-oracle guarantee concerns reusable algebra; full-call time additionally includes boundary scans, state creation, canonicalization and wrapper costs. Native external baselines test the remaining implementation and accuracy gap. Small pilot sample sizes, censored non-completions, and untuned resource settings limit interpretation to the measured implementation and workloads.
\subsection{Independent-process timing replication}
Three fresh-process replicates follow the controlled pilot at the primary surface-code point and the high-noise BB144 point. Each uses the same first 64 inputs with randomized method order. These are additional timing executions, not new independent logical samples. All repeated predictions match the corresponding stored single-pass predictions.
\begin{table*}[t]\centering\scriptsize
\caption{Three independently initialized timing replicates on 64 reused inputs per operating point. RSS is the maximum separately sampled whole-process footprint across replicates.}\label{tab:replication}
\begin{tabular}{llrrrrr}\toprule Point&Backend&Run 1 mean ms&Run 2&Run 3&Pooled p95 ms&Max RSS MiB\\\midrule
BB144&BPHUF&347.53&347.53&348.37&601.85&137.7\\
BB144&BPSW&285.45&288.64&283.37&454.90&137.7\\
BB144&Tesseract&55.98&55.41&59.46&69.05&117.1\\
Surface d5&BPHUF&132.92&130.67&133.22&343.95&112.4\\
Surface d5&BPSW&122.26&119.00&123.61&328.14&111.8\\
Surface d5&Tesseract&117.66&113.71&115.00&246.90&106.8\\
\bottomrule\end{tabular}\end{table*}
The median paired reduction in mean latency is 8.0\% at the surface-code point and 17.9\% at the BB144 point. The corresponding whole-process RSS measurements show little change between oracle modes. Native Tesseract remains a stronger accuracy comparator on these samples; the timing table also exposes dependence on which fixed sample subset is measured.
A separate diagnostic on eight fixed high-noise BB144 inputs records 81,639 rebuilding basis-insertion calls and 3,100 persistent calls, a 26.3-fold reduction in this repeated-column-processing counter. Both diagnostic outputs are identical. The instrumented times are excluded from benchmark latency. Profiling identifies repeated component creation, DSU root queries and eligibility scans as substantial remaining costs; those paths are preserved by the current oracle replacement.



\section{Conclusion}
Seamwise combines fixed-exterior boundary repair, logical-reach analysis, merge-persistent conditional solving, and exact scheduling of disjoint updates. Accepted repairs preserve the full detector record and decrease the prior configuration cost. The persistent state eliminates repeated reconstruction of inherited detector relations, while the exposed-reach certificate characterizes the logical directions represented by the inserted fault support.

The original compact-circuit policy study demonstrates 30.6--34.5\% fewer logical failures than direct cutting in its primary comparison. The new controlled study tests the implementation cost of the persistent backend and compares full-record native baselines on identical noise samples. Exact output preservation is an algorithmic guarantee under canonicalization; a practical speed advantage and competitive logical accuracy require measurement. These results characterize the measured Python reference and circuit sizes. Native optimized implementations, common-budget accuracy--latency frontiers, and complete 2026 comparator coverage remain to be evaluated.

\appendix
\renewcommand{\thesection}{\Alph{section}}
\section{Detailed circuit and execution protocol}\label{app:protocol}
\subsection{Circuit families and observables}
All circuits start from ideal encoded Bell pairs with noiseless references. A perfect final extraction closes the detector record, and final commuting Bell products monitor both logical Pauli sectors. For $B$ blocks encoding $k_q$ logical qubits each, there are $N_{\rm obs}=2Bk_q$ monitored outputs. The Bell-reference construction follows Appendix A.7 of Ref.~\cite{tour}.

Surface-code circuits use rotated $[[d^2,1,d]]$ codes with $d\in\{3,5,7\}$, $B=4$, $d$ noisy extraction rounds, and $p\in\{0.001,0.003,0.006\}$. Each block uses $d^2-1$ extraction ancillas and one noiseless reference in addition to its data qubits. Each noisy round begins with a nearest-neighbor transversal CNOT matching. Matchings alternate, and directions reverse every two rounds. Ancilla resets have $X$-flip noise, ancilla $H$ gates and all CNOT gates have depolarizing noise, and readout bits flip with probability $p$. Idle noise is omitted.

BB72 and BB144 use the $[[72,12,6]]$ and $[[144,12,12]]$ constructions from Ref.~\cite{bbpre}. Four blocks undergo six noisy slices at $p=0.005$ or 0.015. Inter-block CNOTs and commuting stabilizer-product measurements are ideal except for independent readout flips; each slice applies single-qubit depolarization to data. These controls use a two-sector phenomenological noise model. Each monitors 96 Bell observables. Additional eight-block surface cases at $p=0.003$, $d=3,5$, test multiple seams and support the engineering audit. Table~\ref{tab:grid} lists the circuit and sampling parameters. The distance-seven cases have limited statistical power.

Stim~\cite{stim} generates an undecomposed detector error model with no approximation of disjoint errors. We sample the model with a fixed seed and verify $De=s$ and $Le=\ell$ for every shot. Noiseless tests and 54 injected logical $X/Z$ checks validate the Bell-reference observables. No approximate graphlike projection is used to make the correlated hypergraph acceptable to a matching decoder.

\begin{table*}[t]\centering\small
\caption{Circuit parameters and sample counts for the 15 configurations. Each of BB72 and BB144 uses the listed shot counts. Dashes indicate noise settings outside the grid.}\label{tab:grid}
\begin{tabular}{lrrrrrrrr}\toprule
&&&\multicolumn{5}{c}{Shots at circuit noise parameter $p$}\\
Code&Blocks&Noisy rounds&0.001&0.003&0.005&0.006&0.015\\\midrule
Surface $d=3$&4&3&512&512&--&512&--\\
Surface $d=5$&4&5&256&1,024&--&256&--\\
Surface $d=7$&4&7&64&64&--&64&--\\
BB72, BB144&4&6&--&--&128&--&128\\
Surface $d=3$&8&3&--&256&--&--&--\\
Surface $d=5$&8&5&--&64&--&--&--\\\bottomrule
\end{tabular}\end{table*}

\subsection{Decoder stacks and local row selection}
The four engine settings are BF, BP+LSD0, BP+LSD2, and BP+OSD0. BF uses inversion and one-bit growth; LSD2 uses combination-sweep order two, and OSD0 uses order-zero ordered-statistics postprocessing~\cite{lsd,osd,bfdocs,lsddocs,osddocs}. All share 30 min-sum BP iterations, scaling 0.625, parallel message scheduling, and one native CPU thread. We use Stim 1.16.0 and LDPC 2.4.1. The engine grid covers three postprocessing families and two LSD orders. It does not establish optimality among all correlated decoders.

Each engine is crossed with full joint decoding, direct cutting, and radius-one repair. BF additionally tests halos of one and two blocks, radius-two repair, and rank-screened radius-one repair, giving 16 stacks. The core width is two. In four-block models, radius two unfreezes all nonzero detector columns and therefore performs a full-width re-decode filtered against the cut incumbent. In eight-block models, its four-owner-block bands are proper subsets. For a halo window, the local row set includes every touched row with no remaining outside-window variable, together with all core rows. A failed local solve or unresolved residual triggers a full joint call using the same engine; fallback frequency and its full latency are recorded. Seam replacement uses all touched rows, fixed exterior variables, and a $10^{-10}$ acceptance tolerance. Every returned correction is audited against the complete $D$ matrix.

The frozen resource allocation includes LSD order two and OSD order zero. Higher-order OSD is outside the main grid, so the results do not rank boundary repair against that baseline. The protocol records an initial timing pilot as the allocation rationale; its durations are not used as quantitative evidence in the revised comparisons. Matching-based circuit representations~\cite{fast} and other modern BP variants remain relevant external comparisons.

All conditional targets in the study, including warmups, pass the exact column-space membership check. The guard therefore rejects zero native calls in this dataset. Long durations are retained as measured decoding cost and are not classified as evidence of infeasibility.

\subsection{Sampling, statistics, and multiplicity}
We set the configurations, seeds, shot counts, engine settings, and nine planned contrasts per model before initial sampling. During execution, problematic native calls motivated an implementation amendment: an exact binary column-space guard rejects an infeasible local target before calling the native engine. It implements the existing failed-local-solve/fallback rule and leaves feasible inputs unchanged. All configurations were rerun with this guard, whose cost is included in timing. We made the amendment after inspecting partial outcomes. Sample counts, seeds, hypotheses, and engine settings were unchanged. Methods share the same samples and run in seeded random order after three excluded warmups. There is no outcome-dependent stopping. The 15 models contribute 135 planned paired tests to one Holm family~\cite{holm}.

We use 95\% Wilson intervals for marginal failure rates~\cite{wilson}. For a paired contrast, $n_{10}$ counts baseline failures fixed by the comparison method, and $n_{01}$ counts new failures introduced by that method. The effect is $(n_{10}-n_{01})/N$. We use the exact conditional binomial version of the paired McNemar test~\cite{mcnemar} and conservative 95\% paired-effect intervals obtained by subtracting two 97.5\% Clopper--Pearson intervals~\cite{cp}; the union bound supplies simultaneous coverage for the two discordance probabilities. These effect intervals are not adjusted across all contrasts. Appendix~\ref{app:statistics} gives the formulas.

\subsection{Timing and memory estimands}
Main timings include decoding, wrapper operations, consistency checks internal to the policy, and repair-event construction; compilation, sampling, file writes, and the final outcome audit are excluded. Model-level simulation jobs run with up to three concurrent processes on one shared host, each native decoder using one thread; the principal distance-five case was completed with one model job active. These timings describe that execution environment. The separate engineering audit serializes the compared method runs and provides the controlled serial/parallel resource comparison.

The engineering audit runs three fresh-process replicates of each method on 64 common eight-block samples, with seeded randomized method order. It compares event-free serial and two-worker process-parallel repair with identical checks and acceptance arithmetic. Initialization is reported separately; the processes and local solvers remain warm during timed calls. Memory is instrumented in a separate pass over the first eight audited inputs. Per-process OS high-water RSS, covering the preceding timed calls as well, and 2 ms samples of aggregate process-tree RSS and proportional set size (PSS) are retained. Sampled aggregate peaks can miss short transients and uninstrumented inputs; PSS avoids counting shared pages in full for every process.

\section{Additional surface-code and rank results}\label{app:surface}
Figure~\ref{fig:surface} shows the original surface-code sweep at distances three, five, and seven. Within the 15 configurations, radius-one BF repair yields lower observed failure counts than BF cutting in 13 cases and ties in two. All paired contrasts remain in the 135-test correction family. Figure~\ref{fig:rankcurve} traces logical reach as fault columns are added to a repair band. The original figures are retained; the controlled backend study additionally reuses the stored sample archives and records their hashes.

\begin{figure*}[t]\centering\includegraphics[width=\textwidth]{figures/fig5.pdf}
\caption{Surface-code sweep at distances 3, 5, and 7. Methods use the same fault samples within each configuration. Error bars show marginal 95\% Wilson intervals.}\label{fig:surface}\end{figure*}
\begin{figure*}[t]\centering\includegraphics[width=\textwidth]{figures/fig6.pdf}
\caption{Logical reach as the selected fault-column set grows. Columns enter in fixed prior order. The rank calculation uses the circuit matrices without shot-level logical truth.}\label{fig:rankcurve}\end{figure*}

\section{Complete engineering measurements}\label{app:engineering}
The audit serializes the measured method runs. Each method uses three fresh processes on 64 shared samples. Mean and p95 pool those durations; memory entries are maxima of separately instrumented eight-input passes. The main-text speed ratio excludes parallel-worker startup. Table~\ref{tab:engineering} includes both measured distances.
\begin{table*}[t]\centering\small
\caption{Warm-call latency and process-tree memory for the eight-block surface-code circuits. Timing and memory were measured in separate passes.}\label{tab:engineering}
\begin{tabular}{lrrrrrr}\toprule
&\multicolumn{3}{c}{$d=3$}&\multicolumn{3}{c}{$d=5$}\\
Stack&Mean ms&p95 ms&PSS MiB&Mean ms&p95 ms&PSS MiB\\\midrule
BF joint&2.58&5.99&66.6&68.78&149.46&74.6\\
BF cut&2.57&4.86&68.9&35.44&77.20&82.3\\
BF repair1&3.53&6.57&70.3&60.00&137.54&88.1\\
BF parallel1&4.57&8.19&193.0&49.71&122.24&233.1\\
LSD2 joint&2.57&6.73&66.7&54.39&110.84&74.7\\
LSD2 repair1&3.39&6.19&70.0&61.77&155.71&87.7\\
OSD0 joint&2.50&5.26&67.4&51.34&74.21&79.8\\
OSD0 repair1&3.22&6.68&72.3&51.86&73.74&101.6\\\bottomrule
\end{tabular}\end{table*}

\section{Ownership-based initial partition pass}\label{app:initial}
For a core width $q$, process core block sets
\[C_t=\{t,\ldots,\min(t+q-1,B)\},\quad t=1,1+q,\ldots,\ t\le B.\]
The halo width $h$ extends the current owner range to
\begin{equation}W_t=\{j:t\le\own(j)\le\min(t+q+h-1,B)\}.\label{eq:halo}\end{equation}
The inner decoder receives the residual syndrome restricted to a local row set $U_t$. The essential interface requirement is that $U_t$ include every detector row in the core. Additional rows provide context; their exact selection is an implementation choice. After a local candidate is obtained, only variables $K_t=\{j:\own(j)\in C_t\}$ are committed. Halo-owned estimates are discarded. A local candidate must close the core under the residual syndrome; a failure to do so must be flagged rather than silently committed.

\AlgBox{Algorithm A: ownership-based pre-commit overlap}{
\textbf{Input:} $D,s,\own,q,h$, row policy $U_t$, solver $\mathcal A$. Initialize $\hat e=0$, $s^{\rm res}=s$.
\begin{enumerate}
\item For $t=1,1+q,\ldots$, form $C_t,K_t,W_t,U_t$.
\item Compute $z\leftarrow\mathcal A(D_{U_t,W_t},s^{\rm res}_{U_t},p_{W_t})$. If no candidate closes every core row, return FAIL.
\item Set $\hat e_{K_t}\leftarrow z_{K_t}$ and $s^{\rm res}\leftarrow s^{\rm res}+D_{:,K_t}z_{K_t}$. Discard $z_{W_t\setminus K_t}$.
\item Return $\hat e$ only if $D\hat e=s$; otherwise return FAIL.
\end{enumerate}}
\begin{theorem}[Core closure]\label{thm:closure}
Suppose each commitment in Algorithm A closes the detector rows in its core. Then every closed row remains closed after all subsequent commitments. If all cores complete, the returned estimate satisfies $D\hat e=s$.
\end{theorem}
\begin{proof}
Consider a row $i$ in the current core, whose largest block label is $a$. Every column incident on $i$ has owner at most $\beta(i)\le a$. Earlier-owned columns have already been committed, and all remaining incident columns belong to the current core; later halo-owned columns contribute zero to row $i$. Thus a locally consistent candidate closes $i$ when its core variables are committed. Every subsequent committed column has owner greater than $a$ and is zero on row $i$ by Eq.~\eqref{eq:triangular}. Closure is therefore permanent. Induction over the cores closes every row.
\end{proof}
Core closure guarantees syndrome consistency when every core solve succeeds. A committed prefix can satisfy its current rows while making a later core infeasible. Appendix~\ref{app:examples} gives an example and motivates the joint fallback.

\paragraph{Numerical example of local submatrices.}
Return to Eq.~\eqref{eq:toy} and assign detector rows to blocks by $\beta(1)=1$, $\beta(2)=2$. The four column owners are $(1,1,2,1)$. For a one-block first core and one halo block, $C_1=\{1\}$, $K_1=\{1,2,4\}$, and $W_1=\{1,2,3,4\}$. With $U_1=\{1,2\}$, the local matrix $D_{U_1,W_1}$ is the full two-row matrix in Eq.~\eqref{eq:toy}; its solution uses all four columns as context, but only columns 1, 2, 4 are committed. These subscripts always select detector rows first and fault columns second.

The same columns illustrate conditioning and the rank identity. Select $J=\{1,2,4\}$, whose exterior is $\bar J=\{3\}$ and whose touched-row set is $R(J)=\{1,2\}$. Then
\[
D_J=\begin{pmatrix}1&1&1\\0&1&1\end{pmatrix},\quad D_{R(J),\bar J}=\begin{pmatrix}0\\1\end{pmatrix},
\]
\[L_{:,J}=\begin{pmatrix}1&0&1\end{pmatrix},\qquad t_J=\begin{pmatrix}s_1\\s_2+x_3\end{pmatrix}.\]
The exterior contributes a known offset to the second detector equation. For $s=(1,1)^T$ and $x_3=0$, both local candidates $(0,1,0)^T$ and $(0,0,1)^T$ satisfy $D_Jz=t_J$, but their logical actions are zero and one. Their difference $(0,1,1)^T$ lies in $\ker D_J$. Stacking the logical row below the detector rows gives
\[\begin{pmatrix}D_J\\L_{:,J}\end{pmatrix}=\begin{pmatrix}1&1&1\\0&1&1\\1&0&1\end{pmatrix},\quad\rho(J)=3-2=1.\]
Thus the stacked matrix counts constraints on both syndrome and logical action, whereas $D_J$ counts syndrome constraints alone. Their rank difference identifies the one logical bit that the local replacement can change. The acceptance test then chooses whether a proposed replacement lowers the prior-weight cost.

\section{BB matrices and logical basis validation}\label{app:bb}
\subsection{Translations and polynomial matrices}
Let $S_\nu$ be the cyclic shift acting on a unit row vector by $u_aS_\nu=u_{(a+1)\bmod\nu}$. For example,
\[S_3=\begin{pmatrix}0&1&0\\0&0&1\\1&0&0\end{pmatrix}.\]
On a periodic $\ell_x\times\ell_y$ lattice, define $T_x=S_{\ell_x}\otimes I_{\ell_y}$ and $T_y=I_{\ell_x}\otimes S_{\ell_y}$. A row indexed by site $(a,b)$ moves to $(a+1,b)$ under $T_x$ and to $(a,b+1)$ under $T_y$, with coordinates taken modulo the periods. Each binary monomial therefore specifies a translated support. With the code specification of Ref.~\cite{bbpre}, set
\[P=T_x^3+T_y+T_y^2,\qquad Q=T_y^3+T_x+T_x^2,\]
\[H_X=[P\mid Q],\qquad H_Z=[Q^T\mid P^T].\]
Here $P,Q$ each have size $\ell_x\ell_y\times\ell_x\ell_y$; concatenating them creates $n_q=2\ell_x\ell_y$ physical-qubit columns. Each row of $H_X$ gives the support of an $X$ stabilizer, and each row of $H_Z$ that of a $Z$ stabilizer. For BB72, $(\ell_x,\ell_y)=(6,6)$. Number sites within each half by $1+6a+b$. The first row of $P$ has ones at sites 19, 2, 3, and the first row of $Q$ at sites 4, 7, 13. Consequently the first row of $H_X$ specifies $X_2X_3X_{19}X_{40}X_{43}X_{49}$. This is one concrete stabilizer behind the compact polynomial notation. Since translations commute, $PQ=QP$, and
\[H_XH_Z^T=PQ+QP=0.\]
This binary product counts each pair's shared support modulo two, verifying that the $X$ and $Z$ stabilizers commute. The periods $(12,6)$ give BB144.

\subsection{Logical-operator support matrices}
The rows of $G_X$ and $G_Z$ are physical-qubit supports of logical Pauli operators. An $X$ support must be in $\ker H_Z$ to commute with every $Z$ check; adding a row combination of $H_X$ multiplies it by an $X$ stabilizer and preserves its logical action. Thus representatives are selected from $\ker H_Z/\operatorname{row}H_X$, and analogously $Z$ representatives from $\ker H_X/\operatorname{row}H_Z$. Kernels here are identified with their transposed row-vector supports. Canonical pairing requires
\[G_XH_Z^T=0,\qquad G_ZH_X^T=0,\qquad G_XG_Z^T=I_{k_q}.\]
For a small illustration, the $[[4,2,2]]$ CSS code has
\[H_X=H_Z=\begin{pmatrix}1&1&1&1\end{pmatrix},\]
\[G_X=\begin{pmatrix}1&1&0&0\\1&0&1&0\end{pmatrix},\quad G_Z=\begin{pmatrix}1&0&1&0\\1&1&0&0\end{pmatrix}.\]
All logical rows have even weight and hence commute with the checks. The first $X$ row and first $Z$ row overlap on one qubit and anticommute; different logical pairs overlap on an even number and commute. Direct multiplication gives $G_XG_Z^T=I_2$. We use the same quotient-and-pairing procedure for the BB codes.

The distinction from Section III is the matrix axis: $H_X,H_Z,G_X,G_Z$ have physical-qubit columns, whereas $D,L$ have circuit-fault columns. The code-level logical supports define the measured logical observables. Propagating each circuit fault and recording the observable flips then constructs the corresponding column of $L$.

\section{Completion and likelihood examples}\label{app:examples}
\subsection{Completion after a partial cut}
For prefix closure, take
\[D=\begin{pmatrix}1&1&0\\0&1&1\\0&1&0\end{pmatrix},\quad s=(1,1,1)^T=D(0,1,0)^T.\]
Its rows are the equations $x_1+x_2=1$, $x_2+x_3=1$, and $x_2=1$; its columns record which equations a fault toggles. With one row per block, columns one and two are owned by block one, and column three by block two. A first commitment $(x_1,x_2)=(1,0)$ closes row one; $x_3=1$ closes row two. Row three is then unsatisfied with no remaining variable. Earlier rows stay closed, but no completion exists. This is why the reference implementation declares a local failure and includes the same-engine global fallback in its measured cost.

\subsection{Configuration likelihood and logical-class weight}
For independent fault probabilities, logical-class weight is
\[Z_\lambda(s)=\sum_{x:Dx=s,\ Lx=\lambda}\exp[-\cost(x)].\]
The repair objective selects configurations, while the recorded logical outcomes evaluate the complete decoder. For logical degeneracy, take $D=(1\ 1\ 1)$, $L=(1\ 0\ 0)$, $s=1$, and independent probabilities $(.3,.2,.2)$. The detector row checks the parity $x_1+x_2+x_3$, while the logical row records $x_1$. Thus columns two and three have the same detector and logical signatures, and column one has the same detector signature but a different logical action. The syndrome-compatible configurations $100,010,001,111$ have masses $.192,.112,.112,.012$. Replacing $010$ by $100$ improves configuration likelihood, but changes from logical class zero, of total mass $.224$, to class one, of mass $.204$. A lower-cost representative can therefore be in the less probable logical class. This is an algebraic counterexample, not another simulated code experiment.

\section{Paired tests and uncertainty}\label{app:statistics}
For $f$ failures in $N$ independent samples, $\hat P=f/N$ and $z_W=1.9599639845$, the marginal Wilson endpoints are
\[\frac{\hat P+z_W^2/(2N)\ \pm\ z_W\sqrt{\hat P(1-\hat P)/N+z_W^2/(4N^2)}}{1+z_W^2/N}.\]
For paired discordances $n_{10},n_{01}$ and $M_{\rm disc}=n_{10}+n_{01}>0$,
\[p_{\rm exact}=\min\!\left\{1,\ 2\sum_{j=0}^{\min(n_{10},n_{01})}\binom{M_{\rm disc}}j2^{-M_{\rm disc}}\right\}.\]
Set $p_{\rm exact}=1$ for $M_{\rm disc}=0$. Holm adjustment sorts the 135 planned $p$-values and sets
\[\widetilde p_{(i)}=\min\{1,\max_{j\le i}(136-j)p_{(j)}\}.\]
Let $[a_{10},b_{10}]$ and $[a_{01},b_{01}]$ be 97.5\% Clopper--Pearson intervals for the two marginal discordance probabilities. The paired-effect interval $[a_{10}-b_{01},b_{10}-a_{01}]$ has at least 95\% coverage by the union bound, including zero-discordance cases. It is conservative and does not replace the exact paired test.

\section{Evaluation grid}\label{app:grid}
Table~\ref{tab:failgrid} reports the original observed failure counts. R1 and R2 denote repair at radii one and two. Radius-one BF repair lowers the observed failure count relative to direct cutting in 13 of the 15 configurations and ties in the other two. The reductions span both code families. For example, repair reduces failures from 55 to 39 on 1,024 paired shots for the four-block distance-five surface-code circuit at $p=0.003$, and from 44 to 32 on 128 paired shots for BB72 at $p=0.015$. The magnitude of the improvement varies across configurations, so the counts are best read alongside the shot totals.

Increasing the repair radius has mixed effects. Compared with radius-one BF repair, radius-two repair yields fewer failures in five configurations, ties in five, and yields more failures in five. The comparison with joint decoding also depends on the circuit. At $p=0.006$, radius-one BF repair gives 45 failures versus 50 for joint BF in the distance-five circuit, while the distance-seven circuit gives 22 versus 17. These counts describe the observed outcomes. The paired tests account for which individual shots each method fixes or turns into a failure.

\begin{table*}[t]\centering\scriptsize
\caption{Logical failure counts on the complete evaluation grid.}\label{tab:failgrid}
\begin{tabular}{lrrrrrrrrrrrr}\toprule
&&&&\multicolumn{4}{c}{BF}&\multicolumn{2}{c}{LSD2}&\multicolumn{2}{c}{OSD0}\\
Code&$p$&Blocks&$N$&Joint&Cut&R1&R2&Joint&R1&Joint&R1\\\midrule
Surface $d=3$&.001&4&512&10&7&5&8&8&5&10&5\\
Surface $d=3$&.003&4&512&30&29&25&25&28&24&30&25\\
Surface $d=3$&.006&4&512&59&63&54&56&51&51&59&53\\
Surface $d=5$&.001&4&256&1&2&2&1&0&1&1&2\\
Surface $d=5$&.003&4&1024&46&55&39&36&38&33&46&38\\
Surface $d=5$&.006&4&256&50&65&45&47&47&38&50&43\\
Surface $d=7$&.001&4&64&0&1&0&0&0&0&0&0\\
Surface $d=7$&.003&4&64&2&3&0&0&2&0&2&0\\
Surface $d=7$&.006&4&64&17&29&22&17&16&21&17&24\\
BB72&.005&4&128&6&9&5&5&6&4&6&5\\
BB72&.015&4&128&47&44&32&30&47&30&47&32\\
BB144&.005&4&128&0&0&0&0&0&0&0&0\\
BB144&.015&4&128&12&13&10&7&10&6&12&8\\
Surface $d=3$&.003&8&256&26&26&18&21&23&21&26&18\\
Surface $d=5$&.003&8&64&9&9&3&4&8&4&9&3\\\bottomrule
\end{tabular}\end{table*}

\section{Persistent conditional oracle: algorithm and correctness}\label{app:persistent}
This appendix defines the persistent oracle for $A=D_{R(J),J}$ and $t=t_J$. Vectors below use global labels within the local problem, with supports stored only where nonzero. The target $t_S$ equals $t$ on $C(S)$ and is zero elsewhere. At a singleton check $i$, initialize an empty basis, $r_S=t_i e_i$, and $u_S=0$. At a singleton fault vertex, initialize an empty basis and zero residual and witness. A fault column has an initially false activation flag.

\AlgBox{Algorithm B: merge-persistent local oracle}{
\textbf{On merging disjoint components $S_1,S_2$:}
\begin{enumerate}
\item Concatenate their basis lists in decreasing pivot order; set $r=r_{S_1}+r_{S_2}$ and $u=u_{S_1}+u_{S_2}$.
\item Scan fault vertices in the merged component. Activate each not-yet-inserted $j$ whose complete check neighborhood is now internal.
\item For an activated column initialize $(v,c)=(A_{:,j},e_j)$. For each old pivot $p$ in decreasing order, if $v[p]=1$, set $(v,c)\leftarrow(v+b_p,c+c_p)$.
\item If $v\ne0$, let $p=\max\supp(v)$ and insert $(b_p,c_p)=(v,c)$. If $r[p]=1$, set $(r,u)\leftarrow(r+v,u+c)$.
\item If $v=0$, no detector pivot is added. Optionally emit $k_j=c$ for the logical-reach certificate.
\item Return valid exactly when $r=0$; on validity return $u$.
\end{enumerate}}

\begin{lemma}[Disjoint support and monotone eligibility]\label{lem:merge}
Under merge-only growth, previously internal fault columns remain internal. Separate current components have disjoint check-coordinate sets, so their stored detector bases are jointly independent and already in echelon form when ordered by pivot. Their residuals have zeros at all inherited pivots after merging.
\end{lemma}
\begin{proof}
If $f_j$ and all checks adjacent to it lie in $S$, they remain in every later component containing $S$. Hence a column is activated at most once. Each basis vector in $S$ is a combination of internal columns, so its detector support lies in $C(S)$. For distinct components these supports are disjoint. A cross-component linear dependence separates into two independent within-component dependences, each zero. Inherited pivot coordinates are distinct, and each component's residual is zero at its own pivots and outside its check support. Adding residuals preserves all inherited zeros.
\end{proof}

\begin{lemma}[Insertion invariant]\label{lem:insert}
Algorithm B preserves Eq.~\eqref{eq:invariant}. The coefficient attached to a nonzero new detector vector satisfies $Ac=v$, and a dependent insertion emits a true kernel relation.
\end{lemma}
\begin{proof}
Initially $Ac=Ae_j=A_{:,j}=v$. Each XOR with $(b_p,c_p)$ preserves $Ac=v$. Decreasing-pivot reduction clears all old pivots. If $v\ne0$, its leading coordinate is a new pivot and its addition enlarges the span by precisely the new column. If $r[p]=1$, adding $(v,c)$ to $(r,u)$ preserves $r+Au=t_S$ and clears the new pivot without altering old residual pivots. If $r[p]=0$, the invariant already holds. When $v=0$, the column was dependent, the old span already contains it, and $Ac=0$. Its coefficient vector is supported on previously inserted columns and the new column.
\end{proof}

\paragraph{Proof of Theorem~\ref{thm:oracle}.}
Induct over merges and insertions using Lemmas~\ref{lem:merge} and~\ref{lem:insert}. The stored span equals the span of all eligible columns, and $t_S=r_S+Au_S$. A nonzero vector in an echelon span has a nonzero leading pivot: choose the largest pivot among the participating basis vectors. Since $r_S$ is zero at all pivots, it belongs to that span exactly when $r_S=0$. Thus $t_S$ belongs to the eligible column space exactly when the oracle reports valid. On validity, $Au_S=t_S$. Identical validity flags give identical subsequent growth choices for a deterministic scheduler with identical weights and tie handling.\hfill$\square$

\paragraph{Witness choice and canonicalization.}
Two solutions of $Az=t$ differ by a vector in $\ker A$. Their logical images can differ, as Eq.~\eqref{eq:toy} demonstrates. Accordingly, feasibility-oracle equivalence alone cannot transfer a failure rate from one witness policy to another. To reproduce a fixed rebuilding HUF implementation exactly, retain the same terminal clusters and their eligible column sets, then run its specified deterministic Gaussian-elimination rule once on each terminal system. The resulting correction vector is identical, provided all solution-selection details, including free variables and column order, are fixed. A different reliability-ordered optimization postprocessor is not covered merely by naming it canonicalization; its cost must be accounted for separately.

\paragraph{No unrestricted reuse across targets or bands.}
Persistence in the theorem is within one fixed conditional solve. When a new shot or an overlapping repair changes $t_J$, the old residual/witness does not automatically satisfy Eq.~\eqref{eq:invariant}. It must be recomputed or explicitly rebased. Changing $J$ can change the graph and growth trajectory; deleting columns cannot be handled by the merge-only proof. Disjoint components inside one solve have disjoint detector support, whereas distinct Seamwise bands may share detector rows. They therefore run separate oracle states; Theorem~\ref{thm:commute}, rather than direct basis concatenation across bands, justifies their concurrent execution.

\section{Complexity derivations and scope}\label{app:complexity}
\subsection{Representation and graph size}
Fix one call with $m$ nonzero fault columns. Assume maximum detector-row degree $\Delta_c$ and column degree $\Delta_e$ are constants independent of circuit size. Remove zero-degree zero-target rows; any zero-degree nonzero-target row is infeasible. For a retained graph with $n$ checks, $v=n+m$ vertices, and $E=\operatorname{nnz}(A)$ incidences,
\[n\le E\le\Delta_e m,\qquad v=O(m),\quad E=O(m).\]
This is why the relevant graph size includes both checks and mechanisms, yet can be expressed as $O(m)$ under bounded column degree. Transversality alone does not establish this degree assumption for arbitrary propagated circuit faults.

We count word-RAM operations with $O(\log(N_f+N_{\rm det}+2))$-bit indices and bounded-precision scheduler scalars. Store check vectors and coefficient certificates as sorted support-index lists. In a component of size $s$, every such list has at most $s$ entries. Symmetric-difference XOR costs $O(s)$ word operations. The analyzed DSU uses union by size or rank together with path compression; a separate minimum-vertex label supplies deterministic scheduler ties independently of the DSU representative. A full global $N_f$-bit vector per basis element is not the representation underlying the cluster-sensitive theorem. Literal bit-complexity includes an additional index-width factor; arbitrary-precision likelihood arithmetic requires its own cost model.

\subsection{Eliminate each eligible column once}
A component never splits. Lemma~\ref{lem:merge} therefore permits at most $m$ column activations in a call. In a component of size $s\le K$, insertion performs at most $s$ pivot tests and at most $s$ XORs of detector/coefficient vector pairs. Sorted supports make each operation at most $O(s)$, including membership tests if done naively. Thus insertion costs $O(s^2)$, and all insertions cost $O(mK^2)$.

At most $v-1=O(m)$ successful component unions occur. Merging pivot lists, support lists, residuals, and witnesses costs $O(s)$ per union. A simple implementation also scans the merged component's fault vertices and tests all adjacent-check DSU roots to detect newly internal columns. Bounded degree gives $O(s)$ root queries per union. Across the run these costs are $O(mK)$ explicit work and $O(mK\alpha(m))$ amortized DSU work. Already activated columns are skipped for elimination even if scanned again.

\subsection{The weighted scheduler being bounded}
Use nonnegative finite boundary weights and select the smallest unsatisfied component, with a fixed deterministic tie rule. These growth weights may be derived from posterior reliabilities and are distinct from the fixed prior weights used for acceptance. Scan all of the selected component's current boundary incidences, discard incidences whose endpoints are already connected, find the minimum residual growth time, advance its boundary clocks, and merge all newly saturated boundaries. A nonterminal step effects a successful union; since the full local problem is feasible by Eq.~\eqref{eq:incumbenttarget}, an unsatisfied component with no boundary cannot persist. Thus there are $O(m)$ growth steps. Each selected component has at most $K$ vertices and $O(K)$ incidences under bounded degree. Boundary scans and DSU queries total $O(mK\alpha(m))$.

Maintain unsatisfied components in a priority queue keyed by size and deterministic label. Retiring/adding component entries only at a growth selection or merge uses $O(m)$ heap operations overall; lazy entries must likewise be bounded by those events. Queue work is $O(m\log(m+2))$. No priority-queue decrease-key is charged for every scanned incidence: boundary times are computed by scans. A scheduler that performs more updates must include them explicitly. These assumptions define the analyzed implementation and do not assert the same overhead for every weighted-HUF program.

\subsection{Single-call and composed bounds}
Combining algebra and growth yields
\begin{align*}
T_a=O\big(&Im+mK^2+mK\alpha(m)\\
&+m\log(m+2)\big).
\end{align*}
The cross term satisfies $K\alpha(m)\le(K^2+\alpha(m)^2)/2$, and $\alpha(m)^2=O(\log(m+2))$. This gives the summand in Eq.~\eqref{eq:workexact}. A terminal cluster of size $s$ needs $O(s^3)$ work for optional canonical Gaussian elimination. Terminal clusters are disjoint and $s\le K$, so
\begin{equation}\sum_{S\,\rm terminal}|S|^3\le K^2\sum_{S\,\rm terminal}|S|=O(mK^2).\label{eq:canonicalcost}\end{equation}
The per-call basis storage is $O(\sum_{S\,\rm current}|S|^2)\le O(mK)$. The graph, DSU, queue, residuals, and witness supports need no larger asymptotic space. BP with $I$ constant-time message sweeps costs $O(Im)$ on a bounded-degree graph.

Sum the call costs to prove Eq.~\eqref{eq:workexact}, then use $K_a\le K_*$, $m_a\le W$, and $\sum_a m_a=W$ to obtain Eq.~\eqref{eq:work}. A sequential execution retains the global incumbent and one call's state at a time. Concurrent bands instead require the sum of their active local state sizes. Offline computation of $\rho(J)$ and the optional exposed-reach instrumentation are excluded from the core oracle bound and charged below.

\paragraph{Sharper cluster-volume accounting.}
Let $\mathcal T_a$ be the final Tanner components in call $a$, including untouched singletons. Every earlier component lies within one terminal component. Writing $s_C=|C|$, the same charging argument gives the finer bound
\begin{equation}
T_a=O\!\left(Im_a+m_a\log(m_a+2)+\sum_{C\in\mathcal T_a}s_C^3\right).\label{eq:volume}
\end{equation}
Indeed, each terminal component contains at most $s_C$ inserted columns and successful unions, each with at most $O(s_C^2)$ combined algebra/scan cost; the DSU factor is absorbed since $\alpha(s_C)=O(s_C)$. This bound accounts for small clusters individually instead of charging every column at the largest observed $K_a$. It is still an upper bound, not a measured runtime law.

\subsection{Repair work versus complete decoder work}
Let $T_{\rm init}$ be the ownership pass including any initial global fallback, $T_{\rm prep}$ model construction and offline rank computation, and $T_{\rm extra}$ proposal search and wrapper/observable instrumentation beyond persistent HUF and terminal canonicalization. Then
\[T_{\rm total}=T_{\rm prep}+T_{\rm init}+T_{\rm repair}+T_{\rm extra}.\]
No term can be silently removed by calling the scheme local. For a prespecified number of passes and bounded column-visit multiplicity $\mu$, $W\le\mu N_f$. For a single disjoint radius-one batch, $\mu\le1$ after zero columns are fixed. More general overlapping bands require an explicit multiplicity count.

Strict configuration-cost descent ensures finitely many accepted updates, but permits as many as $2^{N_f}-1$ distinct accepted transitions in the absence of stronger structure. Rejected revisits need not terminate at all without a stopping policy. Consequently the work theorem is parameterized by actual or budgeted $W$, not by an unproved polynomial convergence claim. The theorem also says nothing about how often a valid HUF proposal passes the likelihood filter.

\begin{corollary}[All-stage HUF solve work]
Suppose the initial pass and any global fallback also use the analyzed persistent HUF backend. Let $W_{\rm all}$ sum fault-column incidences over initial, fallback, and repair solves, and let $K_{\rm all}$ be their maximum cluster size. If an infeasible local solve terminates when an unsatisfied component has no boundary, then the total HUF solve work is
\[
O\!\left(IW_{\rm all}+W_{\rm all}K_{\rm all}^{2}+W_{\rm all}\log(W_{\rm all}+2)\right).
\]
For core width two, zero halo, one radius-one repair pass, and at most one initial global fallback, $W_{\rm all}\le3N_f$.
\end{corollary}
\begin{proof}
A solve that declares infeasibility makes no more growth events or insertions than the same merge-only analysis allows; the failure branch terminates that call. Sum the per-call bound over all stages. Zero-halo initial cores have disjoint column ownership and visit at most $N_f$ columns in total; a single global fallback visits at most $N_f$; radius-one bands between width-two cores are pairwise disjoint and visit at most $N_f$. Their sum gives the incidence bound. Model preparation, offline ranks, and additional wrapper/observable instrumentation are charged separately.
\end{proof}
The global fallback can determine $K_{\rm all}$, so the incidence bound alone does not imply a small-cluster regime. The measured full-call latency also includes policy bookkeeping and logical-event construction; the solve-work bound does not erase those costs.

\subsection{Conditional small-cluster regime}
For a fixed repair schedule, suppose a run-dependent $K_*$ satisfies
\[\Pr(K_*\ge k)\le\min\{1,CW e^{-ck}\},\qquad c,C>0,\]
with constants independent of size and a deterministic $W$ (or a uniform tail conditional on that $W$). Set $k_0=\lceil c^{-1}\log(CW+2)\rceil$. The tail identity
\[\mathbb E[K_*^2]=\sum_{k\ge1}(2k-1)\Pr(K_*\ge k)\]
gives $O(\log^2(W+2))$: the initial part is $O(k_0^2)$, and the remaining geometric tail is $O(k_0+1)$. Therefore
\[\mathbb E[T_{\rm repair}]=O\!\left(IW+W\log^2(W+2)\right).\]
An exponential cluster tail is an explicit additional hypothesis. A quantum decoding threshold or sparse physical noise by itself is not a proof of that tail for the chosen circuit, posterior weights, and HUF schedule. Existing compact-code observations are insufficient to establish it.

\subsection{A local-circuit criterion for bounded detector degree}
\begin{lemma}[Consecutive-slice detector locality]
Consider identical CSS blocks with uniformly bounded stabilizer weight and qubit--check incidence, bounded-depth stabilizer extraction with bounded-arity gates, and at most one transversal CNOT matching between complete extraction slices. Measured X/Z checks use the same generator labels on matched blocks. Suppose each detector compares a current check measurement with its Clifford-propagated predecessor product, and the check history is reset to the latest record after each complete slice. Then row and column degrees of the detector matrix are bounded independently of block count, code distance and number of slices.
\end{lemma}
\begin{proof}
A matching maps a check to a product on at most two blocks, so a detector uses a constant number of current/preceding check records. A Pauli fault before both compared extraction slices flips those records consistently under the tracked Clifford propagation and cancels in their detector parity. Its noncancelling noisy causal support is therefore confined to a constant number of adjacent bounded-depth extraction/matching circuits. Bounded gate arity and check weight imply that this support contains only a constant number of noisy locations. A constant number of elementary mechanisms per location gives a constant detector-row degree.
Conversely, a fault location has a bounded forward light cone within those adjacent slices. Bounded qubit--check incidence limits the number of affected measurements, and a matching places each previous record in only a bounded number of next detector products. History reset prevents the detector effect from accumulating through arbitrarily many later slices. This bounds column degree as well. Boundary slices only omit predecessors and cannot enlarge the count.
\end{proof}
The generated surface-code family uses the same four extraction-CNOT layers at every distance and resets each check history after measurement. The phenomenological BB controls satisfy the corresponding bounded check-incidence/history argument directly: noisy data Pauli locations and readout flips touch bounded-weight stabilizers under the same matching/history rule. This does not assert a hardware extraction-depth result for BB circuits, which are not benchmarked here. These are the structural conditions used by the criterion; arbitrary deep transversal circuits without the intervening complete extraction/reset structure are not covered. The retained finite models independently record actual row/column degrees in the experiment archive. Logical-action signatures may propagate much farther, so the lemma does not impose a corresponding sparsity bound on $L$.

\subsection{From mechanisms to transversal gate counts}
Let $G$ count logical transversal CNOT instructions, $q(d)$ the data qubits in a participating block, and $n_r$ attributable syndrome rounds per gate. Assume a bounded number of participating blocks per instruction, $O(q(d))$ noisy physical locations per instruction and per syndrome round, $O(1)$ binary mechanisms per location, and initialization/closure volume no larger than the same asymptotic circuit volume. Then
\[N_f=O(\mathcal A),\qquad\mathcal A=G(1+n_r)q(d).\]
If the propagated detector matrices additionally retain uniformly bounded degree and visit multiplicity is at most $\mu$, Eq.~\eqref{eq:work} gives
\begin{equation}
T_{\rm repair}=O\!\left(\mu\mathcal A\,[I+K_*^2+\log(\mu\mathcal A+2)]\right).\label{eq:gate}
\end{equation}
For a surface-code family $q(d)=\Theta(d^2)$ may be substituted. It is not a general identity for all BB or other qLDPC families. A circuit with $G=0$ or significant independent memory volume should be parameterized directly by total noisy location count; the gate-volume relation above applies only when that volume assumption holds.

\subsection{What is improved relative to a repeated solve}
The published HUF analysis gives $O(M^3\log M)$ under its sparse-matrix assumptions~\cite{cain}; a naive dense solve after $O(M)$ merges gives $O(M^4)$. Persistent elimination gives deterministic $O(MK^2+M\log(M+2))$ for the analyzed scheduler, hence $O(M^3)$ in the full-cluster worst case. This removes repeated recomputation; it is not a proof that any physical experiment exhibits a cubic-to-linear runtime crossover. For Seamwise, $M$ is replaced call by call by the local mechanism count $m_a$, and repeated band visits remain charged in $W$.

\section{Logical-reach certificates from dependent insertions}\label{app:reach}
\subsection{Proof of Proposition~\ref{prop:exposed}}
Order all inserted columns by their activation order, breaking simultaneous activations deterministically. A nonzero insertion creates an independent detector pivot; a zero insertion produces a relation $k_j$ containing coefficient one on the new column $j$ and using only previously inserted columns otherwise. Because dependent columns are never added as detector pivots, the coefficients used to eliminate a later column can be chosen on the independent inserted columns alone.

Each dependent $k_j$ has a unique newly introduced dependent-column coordinate. The collection of these relations is linearly independent by triangularity in activation order. If $r=\rank D_{:,F}$, there are exactly $|F|-r$ dependent insertions. Their number equals $\dim\ker D_{:,F}$, so they form a kernel basis. The argument is unaffected by component merges: inherited pivot vectors remain valid, and all activated columns are eventually part of $F$.

Applying the linear map $L$ to that kernel basis yields a spanning set of $L_{:,F}(\ker D_{:,F})$. Zero extension from $F$ to $J$ proves inclusion in $\reach(J)$. Because the smaller space is a subspace of the larger one, equal dimensions imply equal spaces. This proves Eq.~\eqref{eq:active}. Logical images can be dependent or zero even when the detector-kernel relations are independent; a separate logical-space rank reduction is required.\hfill$\square$

\subsection{Component images need not be independent}
For final components $C$, their eligible detector-column supports are disjoint and
\[\ker D_{:,F}=\bigoplus_C\ker D_{:,F_C}.\]
Nevertheless,
\[\reach_{\rm act}(F)=\sum_C L_{:,F_C}(\ker D_{:,F_C})\]
is generally a sum, not a direct sum, in logical-coordinate space. Two components can realize the same logical-output change. Therefore summing componentwise logical ranks can overcount; the logical images must be reduced in their shared output coordinates. This is distinct from Seamwise's disjoint-detector-support invariant.

\subsection{Optional instrumentation cost}
Let $q_a$ be the number of nonzero logical rows of $L_{:,J_a}$ and $r_a^{\ell}\le\min(q_a,m_a)$ the rank of the exposed logical image. For each dependent relation, its certificate has at most $K_a$ fault indices. Dense $q_a$-bit logical signatures can be formed in $O(q_aK_a)$ elementary operations, then reduced against a logical basis in $O(q_a\min(q_a,m_a))$. At most $m_a$ relations are emitted, giving a conservative extra cost
\begin{equation}
T_{\rm reach}=O\!\left(\sum_a m_aq_a[K_a+\min(q_a,m_a)]\right).\label{eq:reachcost}
\end{equation}
Keeping only independent logical images uses $O(q_ar_a^{\ell})$ entries per call, in addition to the persistent detector state and access to $L$. Coefficient witnesses for arbitrary requested logical offsets, if desired, require additional stored relations and are not included in this rank-only memory bound. No constant $q_a$ or bounded logical-signature degree is assumed implicitly.

Offline full-band rank $\rho(J)$ is separate. Straight dense elimination of the stacked $D_J,L_{:,J}$ costs $O((|R(J)|+q_a)m_a\min(|R(J)|+q_a,m_a))$ elementary operations. This can be amortized over repeated shots on a fixed band, but cannot be hidden in a one-shot end-to-end latency claim. An implementation may obtain $\rho(J)$ by processing every band column and collecting dependent logical signatures with the same construction; that offline pass has its own charged work.

\subsection{What full reach does and does not imply}
Let $x$ and an ideal comparison correction $x^\star$ be syndrome-consistent. A replacement supported on $J$ can match the monitored logical output of $x^\star$ exactly when
\[L(x+x^\star)\in\reach(J).\]
This is an existence statement: an admissible difference must lie in $\ker D_J$, and its logical image is precisely an element of $\reach(J)$. If $\reach(J)=L(\ker D)$, every globally attainable logical offset is representable locally. Neither statement ensures that HUF finds it, that its configuration cost is lower, or that it belongs to the most probable logical class. The sampled logical truth is not an input to the repair algorithm.

The active certificate applies to all corrections supported in $F$. A particular terminal-HUF witness policy may explore a still smaller set, because it fixes a solution within each terminal component. Thus $\rho_{\rm act}$ is an upper bound on policy-accessible logical choices; equality with $\rho(J)$ does not establish that the current proposal mechanism searches all of them. This distinction prevents an algebraic reach certificate from being presented as a decoding-accuracy theorem.

\subsection{Reproducible finite algebra checks}
The accompanying \texttt{verify\_integration.py} independently compares persistent validity with rebuilding, checks every returned witness, and compares the streamed logical-image rank against the stacked-rank identity. It covers all nonzero-column $2\times3$ binary matrices, all two-bit targets, all one-row logical maps, and three prescribed incidence orders, giving 2,592 small merge traces; a further 800 seeded traces use matrices with up to ten checks and sixteen fault columns. The run gives 19,721 oracle comparisons, 14,068 valid-witness checks, and 19,721 kernel/logical-rank checks, all without violations. The conditioned-band enumeration verifies 368 feasible replacements across 240 band/incumbent pairs in Eq.~\eqref{eq:toy}, and 1,024 disjoint-update compositions. The three-column example in Section~\ref{sec:work} yields the predicted exposed-rank transition $0\to1$. The generated \texttt{integration\_audit.json} records these counts and the script hash. These finite checks exercise the new implementation and certificate; the proofs above establish the general statements. The script does not sample a quantum circuit, run BP, or benchmark the source experiment counts.

\section{Controlled circuit experiments and reproducibility}\label{app:study}
\subsection{Shared circuits and fixed noise records}
The reproducibility archive contains the recovered original circuit generator, detector matrices, logical-action matrices, fault priors, and stored syndrome/observable samples. Every backend receives the same syndrome rows and is scored against the same full logical-output vector. The primary four-block distance-five case at $p=.006$ requests all 256 stored shots; the other original configurations request the first 64 stored shots, fixed before execution. Sample and matrix SHA-256 values are retained per run. The new scaling grid uses the same generator with independently fixed seeds: $(d,B,T)=(3,16,3),(3,32,3),(3,4,12),(5,8,10),(7,4,14)$, all at $p=.003$, with 64 shots requested per configuration. Here $T$ counts noisy slices, each with the same original transversal matching and extraction rule. A separate fixed-depth sweep adds $d=5,7,9,11$ with four blocks, three noisy slices, $p=.003$, and 64 shots per configuration, using seeds $2026110940+d$; the original distance-three instance anchors that sweep. It extends the maximum code distance to eleven without changing extraction depth along that axis.

\subsection{Exact-oracle ablation and correction selection}
Rebuilding HUF and persistent HUF are two modes of one newly implemented Python reference. They use nonnegative log-prior incidence weights, select the smallest unsatisfied Tanner component, grow every current boundary clock by the minimum residual time, merge all saturated boundaries with deterministic ties, and activate only complete internal fault columns. Both modes perform the same deterministic terminal solve. Their complete Seamwise ownership pass, fallback, radius-one repair, and likelihood filter are identical. The resulting bitstrings and merge-size histories are compared directly. Stored histories cover successfully completed solver calls, including feasible global fallbacks, rather than only accepted seam updates. Infeasible local calls rejected during initialization do not retain their intermediate histories in this prototype; their time remains included in full-call latency. This limits the completeness of the cluster-distribution audit. A zero peak in the BP-assisted variant denotes a call solved by BP before any HUF growth; the supplementary table records this early-convergence fraction. The new implementation is not the published authors' original HUF executable and does not include BP-assisted reliability updates. Consequently its performance is not a claim about an optimized belief-HUF implementation.

An additional BP-assisted residual-HUF ablation uses 30 identical native min-sum BP iterations. A converged BP correction is returned directly. Otherwise both modes subtract the same posterior hard decision from the syndrome, use absolute posterior log-odds as nonnegative growth weights, order terminal columns by that same reliability, decode the residual, and restore the hard decision. This is an explicitly specified variant, not a claim of matching every detail of the original belief-HUF implementation. Thirty-two shared circuit syndromes additionally verify bitwise agreement and syndrome consistency before the full paired pilot.

The reference uses global Python integer bitsets, explicit member scans, and diagnostic histories. It also uses minimum-label DSU representatives and re-sorts pivot keys, rather than the analyzed union-by-rank DSU and maintained sorted pivot lists. It is an auditable prototype rather than the support-local representation analyzed in Theorem~\ref{thm:work}. Wall-clock measurements therefore test the full implementation, not merely the asymptotic number of algebra operations. One hundred independently seeded feasible small Tanner graphs additionally give exact agreement of both correction vectors and complete growth histories between modes.

\subsection{External methods and applicability}
Tesseract~\cite{tesseract} is run through its native Python/C++ package on the identical undecomposed detector model, with detector beam five and priority-queue limit 200,000. It returns predictions for all monitored logical observables. Native BP+LSD2 and BP+OSD2 use the same recorded matrices and priors, 30 min-sum BP iterations, scale .625, and combination-sweep order two. The former is a 2025 clustering baseline, whereas the latter increases the OSD search order beyond the original order-zero grid. External methods decode the full detector record; HUF modes run the complete repair policy. These distinct scopes are reported explicitly rather than relabeled as equal computational workloads.

Relay-BP~\cite{relay} was also selected, but its native package could not be built in this environment because a compatible wheel and Rust toolchain were unavailable. It is recorded as unavailable, not as a failed decoder. The 2026 GARI method~\cite{gari}, ghost protocol~\cite{ghost}, and logical-product approach~\cite{fast} require verified circuit or observable adapters; they have not been benchmarked here. Merely citing these methods does not establish a performance comparison against them. The current study therefore does not complete the desired 2026 baseline coverage.

\subsection{Timing, memory, and compute-budget amendments}
Every method is run in a fresh process, sequentially. Model/solver setup is recorded separately; one warmup is excluded and every subsequent full decode duration is retained. A separate pass over the first four inputs polls whole-process RSS at 1 ms intervals, so memory polling is absent from latency measurements. RSS includes imported libraries, detector models, stored samples, and retained diagnostics; it is not a backend-only memory estimate. Peaks shorter than the polling period may be missed. Native/Python and startup-cost differences must be considered when interpreting these implementation comparisons.

The first low-noise BB144 runs had a 180-second process budget. Following slow high-noise runs, newly started children were capped at 45 seconds including startup. This resource amendment did not use their logical outcomes. Before the additional BP-assisted runs, 180 seconds was preallocated to each primary 256-shot BP-oracle variant; other new runs retain the 45-second cap. Completed fixed-sample runs are summarized with marginal 95\% Wilson intervals; incomplete runs and partial shot journals are retained separately. Because a time limit can censor difficult syndromes preferentially, completed prefixes from interrupted runs are not used as unbiased logical-error-rate evidence. The pilot uses one timing pass per backend and has neither three independent timing replicates nor controlled CPU affinity. Its p99 estimates and small-sample logical rates are descriptive measurements, not tail-latency or threshold guarantees.

The archive contains full per-shot durations, predictions, correction bitstrings, cluster histories, input hashes, source hashes, environment versions, successful summaries, and non-completion logs. The extension measures scale without assuming bounded detector degree: actual maximum row and column degrees are recorded for every completed run. Their finite values do not prove a uniform bound over a growing circuit family.


\subsection{Failure metrics across growing block counts}
The primary error metric remains failure of any monitored logical output, matching the original study. Increasing block count increases the number of monitored outputs, so those whole-circuit probabilities are not fixed-observable threshold comparisons. The supplementary CSV also records the mean logical-observable bit-error rate and monitored-output count. Correlations between observables prevent interpreting that mean as an independent-sample rate without further analysis.

The process budget includes setup, warmup, the timed sample pass, and the separate memory pass. A missing final summary can therefore result from interruption during memory instrumentation even after some or all timed inputs completed; it is not a declaration of decoding infeasibility. All per-shot journals are retained for that distinction.


\subsection{Timing follow-up and structural-work diagnostic}
After the pilot, the primary surface-code point and the high-noise BB144 point were selected for three independent-process timing repetitions, using 64 fixed inputs and randomized method order with seed 774411. The latter point illustrates the large-cluster regime observed during the pilot; this follow-up is not a preregistered survey of all configurations. The 18 repeated runs add timing replication without increasing independent logical sample counts. A separate cProfile diagnostic uses eight fixed BB144 inputs and is excluded from latency summaries. Its insertion-invocation counter quantifies repeated column processing, while its instrumented durations serve only to locate remaining implementation costs.


\begin{thebibliography}{99}\small
\bibitem{cain} M. Cain et al., ``Correlated decoding of logical algorithms with transversal gates,'' \emph{Physical Review Letters} \textbf{133}, 240602 (2024); arXiv:2403.03272v2, Supplementary Sec. 1.2. \url{https://arxiv.org/abs/2403.03272}.
\bibitem{fast} M. Cain et al., ``Fast correlated decoding of transversal logical algorithms,'' arXiv:2505.13587 (2025). \url{https://arxiv.org/abs/2505.13587}.
\bibitem{window} L. Skoric, D. E. Browne, K. M. Barnes, N. I. Gillespie, and E. T. Campbell, ``Parallel window decoding enables scalable fault tolerant quantum computation,'' \emph{Nature Communications} \textbf{14}, 7040 (2023).
\bibitem{bb} S. Bravyi, A. W. Cross, J. M. Gambetta, D. Maslov, P. Rall, and T. J. Yoder, ``High-threshold and low-overhead fault-tolerant quantum memory,'' \emph{Nature} \textbf{627}, 778--782 (2024).
\bibitem{bbpre} S. Bravyi et al., ``High-threshold and low-overhead fault-tolerant quantum memory,'' arXiv:2308.07915v2 (2024), Sec. 4 and Table 3. \url{https://arxiv.org/abs/2308.07915v2}.
\bibitem{tour} T. J. Yoder et al., ``Tour de gross: A modular quantum computer based on bivariate bicycle codes,'' arXiv:2506.03094v1 (2025). \url{https://arxiv.org/abs/2506.03094v1}.
\bibitem{lsd} T. Hillmann, L. Berent, A. O. Quintavalle, J. Eisert, R. Wille, and J. Roffe, ``Localized statistics decoding for quantum low-density parity-check codes,'' \emph{Nature Communications} \textbf{16}, 8214 (2025), doi:10.1038/s41467-025-63214-7. \url{https://arxiv.org/abs/2406.18655}.
\bibitem{osd} P. Panteleev and G. Kalachev, ``Degenerate quantum LDPC codes with good finite length performance,'' \emph{Quantum} \textbf{5}, 585 (2021).
\bibitem{stim} C. Gidney, ``Stim: a fast stabilizer circuit simulator,'' \emph{Quantum} \textbf{5}, 497 (2021).
\bibitem{bfdocs} LDPC project, ``Belief find decoder,'' software documentation, original study access date 4 November 2026. \url{https://software.roffe.eu/ldpc/ldpc/belief_find_decoder.html}.
\bibitem{lsddocs} LDPC project, ``BP+LSD decoder,'' API documentation, original study access date 4 November 2026. \url{https://software.roffe.eu/ldpc/ldpc/bplsd_decoder.html}.
\bibitem{osddocs} LDPC project, ``BP+OSD decoder,'' software documentation, original study access date 4 November 2026. \url{https://software.roffe.eu/ldpc/ldpc/bposd_decoder.html}.
\bibitem{mcnemar} Q. McNemar, ``Note on the sampling error of the difference between correlated proportions or percentages,'' \emph{Psychometrika} \textbf{12}, 153--157 (1947).
\bibitem{holm} S. Holm, ``A simple sequentially rejective multiple test procedure,'' \emph{Scandinavian Journal of Statistics} \textbf{6}, 65--70 (1979).
\bibitem{wilson} E. B. Wilson, ``Probable inference, the law of succession, and statistical inference,'' \emph{Journal of the American Statistical Association} \textbf{22}, 209--212 (1927).
\bibitem{cp} C. J. Clopper and E. S. Pearson, ``The use of confidence or fiducial limits illustrated in the case of the binomial,'' \emph{Biometrika} \textbf{26}, 404--413 (1934).
\bibitem{wiedemann} D. Wiedemann, ``Solving sparse linear equations over finite fields,'' \emph{IEEE Transactions on Information Theory} \textbf{32}, 54--62 (1986).
\bibitem{tesseract} L. Aghababaie Beni, O. Higgott, and N. Shutty, ``Tesseract: A Search-Based Decoder for Quantum Error Correction,'' arXiv:2503.10988 (2025). \url{https://arxiv.org/abs/2503.10988}.
\bibitem{relay} T. M\"uller et al., ``Improved belief propagation is sufficient for real-time decoding of quantum memory,'' arXiv:2506.01779 (2025). \url{https://arxiv.org/abs/2506.01779}.
\bibitem{gari} A. S. Maan, F. M. Garcia Herrero, A. Paler, and V. Savin, ``Decoding correlated errors in quantum LDPC codes,'' \emph{Nature Communications} \textbf{17}, 3965 (2026), doi:10.1038/s41467-026-70556-3.
\bibitem{ghost} ``Scalable Decoding Protocols for Fast Transversal Logic in the Surface Code,'' \emph{PRX Quantum} \textbf{7}, 010320 (2026), doi:10.1103/nx6p-hjqy.
\end{thebibliography}

\end{document}
